% 本文件由 LaTeX 排版 Agent 根据有效 translation.json 自动生成。
% author=Clement of Alexandria; work_id=0208; title=劝勉外邦人

\chapter{劝勉外邦人}

% structure: p0001 source=h1 target=omitted reason=page-title-already-rendered-by-parent

\Signature{\Emph{\Emph{克肋孟·亚历山大里亚}}}\par

\Emph{第一章} 劝勉放弃偶像崇拜的不敬神秘，转而崇拜圣言和圣父天主。\newline{} \Emph{第二章} 外邦神秘和关于诸神生与死之寓言的荒谬与不敬。\newline{} \Emph{第三章} 向诸神祭献的残忍。\newline{} \Emph{第四章} 用来崇拜诸神之像的荒谬与可耻。\newline{} \Emph{第五章} 哲学家们关于天主的观点。\newline{} \Emph{第六章} 哲学家们因神感有时触及真理。\newline{} \Emph{第七章} 诗人亦为真理作证。\newline{} \Emph{第八章} 真道应在先知书中寻求。\newline{} \Emph{第九章} \Quote{藐视或忽视天主恩宠召唤的人犯了大罪。}\newline{} \Emph{第十章} 对外邦人反对说放弃祖先习俗不当的答复。\newline{} \Emph{第十一章} 藉基督的来临，人类获得何等大的恩惠。\newline{} \Emph{第十二章} 劝勉放弃旧时的错误，而听从基督的训诲。\par

\SourceRecord{Translated by William Wilson.；Ante-Nicene Fathers；Vol. 2.；Edited by Alexander Roberts, James Donaldson, and A. Cleveland Coxe.；Buffalo, NY: Christian Literature Publishing Co.,；1885.；<http://www.newadvent.org/fathers/0208.htm>.}

% structure: p0001 source=h1 target=section reason=page-title

\section{劝告异教徒（第一章）}

\Emph{\Person{Clement of Alexandria}}\par

% structure: p0003 source=h2 target=subsection reason=relative-heading-rank; base=h2

\subsection{第一章  劝告异教徒舍弃偶像崇拜的亵渎神秘礼，而归向天主圣言和天主圣父的朝拜}

\Place{特拜}的\Person{安菲翁}和\Place{美提姆纳}的\Person{阿里翁}都是吟游诗人，皆在故事中闻名。至今在希腊人的合唱中仍被歌颂：一位以音乐之力诱鱼，另一位以音乐之力筑墙围困特拜。另一位色雷斯人，精于其艺（他也是希腊传说的主题），仅凭歌声驯服野兽；并用音乐移植橡树。我还可以告诉你们另一位他们的兄弟——神话和歌手的主题——洛克里的\Person{Eunomos}和皮提亚的蟋蟀。庄严的希腊集会曾在皮托聚集，以庆祝皮提亚之蛇的死亡，当时\Person{Eunomos}吟唱爬虫的墓志铭。他的颂歌是赞美蛇的赞歌，还是挽歌，我无法确定。但这是一场比赛，\Person{Eunomos}在夏天弹奏竖琴：正是蟋蟀被太阳晒暖，在山丘的叶子下啁啾的时候；但它们不是在为那条死去的龙歌唱，而是为全智的天主歌唱——一首不拘泥于规则的曲调，胜过\Person{Eunomos}的韵律。洛克里人断了一根弦。蟋蟀跳到琴颈上，在琴上如树枝般歌唱；而吟游诗人顺应蟋蟀的歌声，弥补了缺弦的不足。因此，蟋蟀被\Person{Eunomos}的歌声吸引，正如寓言所描绘，据此在皮托立起一尊\Person{Eunomos}持竖琴的铜像，以及洛克里人在比赛中的盟友。但它自行飞向竖琴，自行歌唱，被希腊人视为一位音乐表演者。\par

请问，你们如何相信虚妄的寓言，以为动物会被音乐迷惑？而真理的独一光辉，似乎在你们看来却乔装打扮，并以不信的眼光看待。因此，\Place{Cithæron}、\Place{Helicon}、奥德里赛的山脉，以及色雷斯人的入教仪式——欺骗的神秘礼——被奉为神圣，并在赞美诗中歌颂。我本人为这些悲剧题材的灾祸感到痛心，尽管只是神话；而你们却将苦难的记录编成戏剧作品。\par

至于那些如今已全然陶醉的戏剧和狂乱诗人，让我们用常春藤冠冕；他们完全疯狂，以酒神的方式，与萨提尔、狂乱的乌合之众以及其余鬼怪之群为伍，让我们把他们禁闭在现已过时的\Place{Cithæron}和\Place{Helicon}。\par

但让我们从天上带来真理，连同一切光明的智慧，以及神圣的先知诗班，下到天主的圣山；让真理将光芒射至最遥远之处，向那些陷入黑暗者四射光芒，并拯救人脱离幻觉，伸出其极其坚强的右手——即智慧——以成就他们的救恩。他们举起双眼，向上观看，要放弃\Place{Helicon}和\Place{Cithæron}，而居住在\Place{熙雍}。\Quote{因为法律将从熙雍发出，上主的话将从耶路撒冷而来，} \BibleRef{依 2:3}——天上的圣言，真正的运动员，在全宇宙的剧场上加冕。我的\Person{Eunomos}所唱的，不是\Person{Terpander}的节奏，也不是\Person{Capito}的节奏，不是弗里吉亚、吕底亚或多利亚的节奏，而是那新的、属神之名的新和音的永恒节奏——新的肋未之歌。\par

\Quote{平息痛苦，平息愤怒，使人忘却一切灾祸。}\par

甜美而真实的是那与这首曲调相融的劝服之咒语。\par

因此，在我看来，那个色雷斯\Person{Orpheus}、那个特拜人、那个美提姆纳人——他们虽为人，却名不副实——似乎是骗子，他们以诗歌为幌子败坏人类生活，被巧诈的巫术精神所占据以达到毁灭的目的，在狂欢中歌颂罪行，把人类的苦难当作宗教崇拜的材料，最先引诱人去敬拜偶像；甚至以木石——即雕像和偶像——建立各国的愚蠢，用他们的歌谣和咒语，使那些原本在天下自由生活的公民遭受极端奴役的枷锁。但我的歌并非如此，它到来是要迅速解除那暴虐魔鬼的苦涩奴役；带领我们回到虔诚的温和可爱之轭，召回那些曾被抛倒在地的人回到天堂。唯有它驯服了人，那最顽固的动物：其中轻浮者如同空中的飞鸟，欺骗者如同爬虫，易怒者如同狮子，纵欲者如同猪，贪婪者如同狼。愚蠢的人是木石，甚至比石头更愚昧的是沉浸于无知的人。作为我们的证人，让我们引用与真理相符的先知之声，为那些被无知和愚昧压碎的人哀叹：\BibleQuote{因为天主能从这些石头给亚巴郎兴起子孙来；} \BibleRef{玛 3:9} 祂怜悯他们极大的无知和顽梗的心，他们对真理石化，从那些石头中——即信赖石头的万民中——兴起了虔诚的种子，敏于德行。再次，一些恶毒虚谎的伪善者，曾图谋不义，祂一度称他们为\Quote{毒蛇的种类}。\BibleRef{玛 3:7} 但如果那些毒蛇中有一条愿意悔改，并跟随圣言，他就成为天主的人。\par

他比喻地将其他人称为披着羊皮的狼，意指那些以人形出现的贪婪怪兽。因此，这首天上的歌将所有这样最凶猛的野兽，所有这样顽石般的生命，都化为温顺的人。\BibleQuote{因为连我们自己也曾经是昏聩的、悖逆的、受迷惑的，服侍各样的私欲和逸乐，生活在恶毒和嫉妒之中，是可憎的，彼此相恨。}宗徒的圣经这样说：\BibleQuote{但当我们的救主天主的良善和施爱于人显示出来时，祂救了我们，并不是由于我们所行的义德，而是按照祂的怜悯。}\BibleRef{铎 3:3}请看\OriginalTerm{新歌}{New Song}的威力！它用石头造人，用野兽造人。此外，那些像死人一样、没有分沾真正生命的人，只因成为这歌的听众就复生了。它还将宇宙组成和谐的秩序，将元素的不协和调成协调的安排，使整个世界成为和谐。它释放了流动的海洋，又阻止它侵占陆地。它又安定曾处于动荡状态的大地，将海定为它的边界。它用大气缓和了火的猛烈，正如多利亚调式与吕底亚调式相混合；又用火的拥抱调节了空气的严寒，和谐地安排了这些宇宙的极端音调。这永不消逝的曲调——万有的支持和全体的和谐——从中心到圆周，又从极端到中心，协调了这万有的架构，不是按照色雷斯的音乐（那就像\Person{犹巴耳}发明的），而是按照天父的旨意，这旨意燃起了\Person{达味}的热情。那出自\Person{达味}、却又先于\Person{达味}的，\OriginalTerm{天主圣言}{the Word of God}，藐视里拉琴和竖琴这些无生命的乐器，并藉\OriginalTerm{圣神}{the Holy Spirit}调谐了宇宙，特别是人——由身体和灵魂构成，是一个小宇宙——在这多音调的乐器上向天主奏乐；并向这乐器——我是指人——唱出相合的歌：\Quote{因为你是我的竖琴、笛子和圣殿。}——竖琴为和谐，笛子因圣神，圣殿因圣言；这样，第一样发出声音，第二样呼吸，第三样容纳主。达味王，我们刚才提到的弹竖琴者，他劝勉真理，劝人远离偶像，他远非在歌中颂扬魔鬼，实际上魔鬼被他的音乐驱走。因此，当\Person{撒乌耳}被魔鬼折磨时，达味只凭弹奏就治愈了他。主按自己的肖像造了人，成为一件美好的、有呼吸的音乐乐器。而那\OriginalTerm{超世智慧}{supramundane Wisdom}、\OriginalTerm{天上的圣言}{celestial Word}自己，也确实是天主全和谐、悦耳、神圣的乐器。那么，这乐器——天主的圣言、主、新歌——渴望什么呢？开盲人的眼，聋子的耳，引瘸子或迷途者进入义德，向愚者显示天主，止息朽坏，战胜死亡，使悖逆的儿女与父亲和好。天主的乐器爱人。主怜悯、教导、劝勉、告诫、拯救、庇护，并因祂的慷慨将天国作为学习的赏报应许给我们；祂所得到的唯一利益就是我们得救。因为邪恶以人的毁灭为养料，而真理如同蜜蜂，不伤害任何东西，只以人的得救为乐。\par

那么，你有了天主的应许；你有了祂的爱：成为祂恩宠的分受者。不要以为\OriginalTerm{救恩之歌}{song of salvation}是新的，像器皿或房屋是新的那样。因为\Quote{在晨星之前它就存在了；}以及\BibleQuote{在起初已有圣言，圣言与天主同在，圣言就是天主。}\BibleRef{若 1:1}错误似乎是古老的，但真理似乎是新事物。\par

那么，是否弗里吉亚人因寓言中的山羊而被证明是最古老的民族？或者另一方面，阿卡迪亚人被那些描述他们比月亮还古老的诗人所证明？又或者，埃及人被那些梦想这片土地首先生出诸神和人类的人所证明？然而，至少这些民族没有一个在世界之前存在。但在奠定世界之前，我们就存在着，我们因注定要在祂内，在祂眼中预先存在——我们\OriginalTerm{天主圣言}{the Word of God}的理性受造物，我们因此从起初就开始算起；因为\BibleQuote{在起初已有圣言。}由于圣言从起初就存在，祂曾是并永远是万有的神圣根源；但由于祂如今取了\OriginalTerm{基督}{Christ}的名字，这名字古时就已被祝圣，并满有能力，祂被我称为\OriginalTerm{新歌}{New Song}。那么，这圣言、基督，既是我们在起初存在的缘由（因为祂在天主内），又是我们福祉的缘由，这同一圣言如今显现为人，祂唯独是两者，既是天主又是人——我们一切福佑的作者；藉着祂，我们被教导善生，被送上通往永生的道路。因为，按照那位蒙启示的主的宗徒所说：\BibleQuote{天主的救恩恩宠已显示给所有人，教导我们弃绝不虔敬和世俗的私欲，在今世生活得清醒、正义和虔敬，期待那蒙福的盼望，以及我们伟大的天主和救主耶稣基督荣耀的显现。}\BibleRef{铎 2:11}\par

这是新歌，是那在起初、在万有之先就已存在的圣言的显现。那先前存在的救主，在近日显现了。那在真实者之内的那一位显现了；因为曾\BibleQuote{与天主同在}的圣言，万物藉祂造成，如今作为我们的导师显现了。那在起初作为创造者创造我们时赋予我们生命的圣言，在作为导师显现时教导我们善度生活；为使祂作为天主，日后引领我们进入永不会终结的生命。祂并非现在才首次怜悯我们的错谬；而是从起初，从万有之先就怜悯我们。但如今，在祂显现时，我们已然失丧，祂完成了我们的救恩。因为那邪恶的爬行怪物，以其蛊惑，至今仍奴役和折磨人类；在我看来，它施以如此野蛮的报复，如同那些把俘虏与尸体捆绑在一起直到一同腐烂的人所说的一样。因此，这个邪恶的暴君和毒蛇，用可悲的迷信锁链，将凡能从出生时就拉拢到他那边的人，捆绑在石头、木块、偶像以及诸如此类的神像上，确实可以说，他把活人与那些死偶像一同埋葬，直到二者一同朽坏。\par

因此（因为诱惑者是同一位），那起初将厄娃引向死亡的，如今也将其余人类引向死亡。我们的盟友和帮助者也是同一位——主，祂从起初便藉着预言赐予启示，如今却明确地呼唤归向救恩。因此，我们当顺从宗徒的训诫，逃避\BibleQuote{空中掌权者的首领，即现今在悖逆之子身上运行的灵}，\BibleRef{弗 2:2}并奔向那现在呼唤我们归向救恩的主救主，祂一向如此，正如祂在埃及和旷野藉着神迹奇事，藉着荆棘和云柱所行的，那云柱因天主圣爱的眷顾，像侍女一样陪伴着希伯来人。藉由这些所激起的敬畏，祂对硬心的人说话；同时，藉着满有智慧的梅瑟、热爱真理的依撒意亚以及全体先知合唱团，以一种更诉诸理性的方式，祂使有耳能听的人转向圣言。有时祂责备，有时祂威胁。有些人祂为之哀叹，另一些人祂用歌唱之声呼唤，正如一位良医，对某些病人用膏药，对某些用按摩，对某些用热敷；有时用柳叶刀切开，有时烧灼，有时截肢，为的是尽可能治好病人患病的部位或肢体。救主有许多音调，有许多拯救人的方法；藉着威胁祂劝诫，藉着责备祂使人悔改，藉着哀叹祂怜悯，藉着歌唱之声祂鼓舞。祂藉燃烧的荆棘说话，因为当时的人需要神迹奇事。\par

祂以火使人敬畏，当祂使火焰从云柱中喷出时——这同时是恩宠与畏惧的标记：你若顺从，就有光明；你若悖逆，就有烈火；但既然人性比云柱或荆棘更高贵，在它们之后，先知们发出了声音——主亲自在依撒意亚、在厄里亚中说话，——亲自藉先知的口说话。但若你不相信先知，反而以为那些人和火都是神话，主将亲自对你说话，\BibleQuote{祂虽具有天主的形体，却没有以自己与天主同等为应得之物，反而空虚自己，} \BibleRef{斐 2:6}——祂，慈悲的天主，竭力拯救人类。如今圣言亲自清晰地向你说话，羞辱你的不信；是的，我说，天主圣言成为人，为使你从人学习人如何能成为神。我的朋友们，这岂非荒谬吗？天主不断地劝勉我们行善，我们却拒绝祂的仁慈，弃绝救恩？\par

难道若翰不也邀请人归向救恩，他岂不完全是劝诫的声音吗？那么，让我们问他：\Quote{你是何人，从何而来？}他必不说他是厄里亚。他将否认自己是基督，却会宣称自己是\BibleQuote{旷野中呼喊的声音。}那么，若翰是谁？\BibleRef{若 1:23}简言之，我们可以说：\Quote{圣言在旷野中恳求呼喊的声音。}你呼喊什么，声音？也告诉我们：\BibleQuote{修直主的道路。}\BibleRef{依 40:3}若翰是前驱，那声音是圣言的先锋；一个邀请的声音，预备救恩——一个催促人奔向天堂基业的声音，藉此，不妊和荒芜的不再无子女。这丰饶是天使的声音所预言的；这声音也是主的先锋，向不妊的妇人传报喜讯，如同若翰向旷野所做的一样。因此，藉着这圣言的声音，不妊的妇人生育子女，旷野结出果实。那预告主来临的两个声音——天使的声音和若翰的声音——据我看来，暗示了为我们预备的救恩：当这圣言显现时，我们将收获这丰饶的果实，即永生。圣经藉着将这两个声音指向同一件事，清楚地说明了这一切：\BibleQuote{不妊的，请听！没有经历产痛的，请发声欢呼！因为荒芜者的子女比有夫之妇的子女更多。}\BibleRef{依 54:1}\par

天使向我们报告了一个丈夫的喜讯。若翰恳求我们认识那位农夫，寻找那位丈夫。因为这使不生育的妇人有丈夫的，并使荒野有农夫的，乃是同一位——祂以天主的德能充满了不生育的妇人和荒野。因为，尽管那位高贵治理之母有许多儿女，但那曾一度多子的希伯来妇人却因不信而成了无子的：不生育的妇人接受了丈夫，荒野接受了农夫；于是二者都因圣言而成为母亲，一者结果实，一者生信者。但对于不信者，不生育的妇人和荒野仍被保留着。故此，圣言的前驱若翰劝勉人们，要为迎接天主的基督的到来做好准备。匝加利亚的哑巴正是这事的预表，它等待基督的前驱结果子，好让作为真理之光的圣言，藉成为福音，打破先知奥秘沉默的预言谜语。但如果你真正渴望看见天主，就当为自己取用配得上祂的净化之物——不是月桂树叶与羊毛和紫色编织的冠冕，而是以公义为冠冕缠绕你的额，以节制的枝叶为环，诚挚地寻求基督。\Quote{因为我是}祂说，\Quote{门}\BibleRef{若 10:9}，我们这些渴望认识天主的人必须找到这扇门，好让祂为我们敞开天堂的大门。因为圣言的门是理智的，由信德的钥匙开启。除了圣子，以及圣子愿意启示的人，没有人认识天主。\BibleRef{玛 11:27} 我深知那开了这扇迄今关闭之门的一位，必将启示里面的东西；祂将显示那些我们若不藉着基督进入便无法知道的事，因为唯独藉着祂，天主才被看见。\par

\SourceRecord{Translated by William Wilson.；Ante-Nicene Fathers；Vol. 2.；Edited by Alexander Roberts, James Donaldson, and A. Cleveland Coxe.；Buffalo, NY: Christian Literature Publishing Co.,；1885.；<http://www.newadvent.org/fathers/020801.htm>.}

% structure: p0001 source=h1 target=section reason=page-title

\section{劝勉外邦人（第二章）}

\Signature{亚历山大的克肋孟著}\par

% structure: p0003 source=h2 target=subsection reason=relative-heading-rank; base=h2

\subsection{第二章  外邦奥秘和关于诸神生死寓言的荒谬与不虔}

那么，不要过于好奇地探究不虔的圣所、充满怪异的洞穴之口、特斯普罗提亚的铜锅、基尔拉的三脚祭台、或多铎那的铜器。曾被视为神圣、位于沙漠中的格兰德里昂及其连同橡树本身已归于衰朽的神谕，被弃于古老寓言之域。卡斯塔利亚泉已沉寂，科罗丰的另一泉也如此；同样，所有其他占卜之泉都已枯竭，被剥去虚夸，尽管为时已晚，却连同其虚构传说一同显现为干涸。也向我们数算那些其他种类的占卜——或不如说是疯狂——毫无用处的神谕：克拉里亚、皮提亚、狄底玛、安菲亚劳斯、阿波罗、安非洛科斯；你若愿意，还可加上解释奇事者、占卜官和释梦者。再带来并放置于皮提亚旁边那些用面粉占卜者、用大麦占卜者，以及仍被许多人尊崇的腹语者。让埃及人的秘密圣所和埃特鲁里亚人的招魂术归于黑暗。这些确实全是不信者疯狂的把戏。山羊也成了这占卜技艺的同谋，被训练来行预言；乌鸦被人教导向人发出神谕般的回应。\par

再者，若我谈论诸奥秘，我并非如他们所说阿尔基比亚德那样加以嘲弄，而是要借真理之言彻底揭露其中隐藏的巫术；并且要将那些作为神秘仪式对象的你们的所谓诸神，仿佛在生命舞台上展示给真理的观众。酒神狂女们举行狂欢以尊崇疯狂的狄俄尼索斯，通过食生肉庆祝其神圣的狂乱，并分配宰杀牺牲的肢体，头戴蛇冠，尖叫着那名为厄娃的名字——因她，谬误进入世界。酒神狂欢的象征是一条被祝圣的蛇。此外，按照希伯来语的严格解释，赫维亚这个含气音的名字意指一条母蛇。\par

得默忒耳和普洛塞耳皮娜成了一出神秘剧的女主角；厄琉西斯以火炬游行庆祝她们的流浪、被掳和哀伤。我以为orgies和mysteries的语源应分别追溯：前者来自得默忒耳对宙斯的\OriginalTerm{愤怒}{\GreekText{ὀργή}}，后者来自与狄俄尼索斯相关的邪恶（\OriginalTerm{玷污}{\GreekText{μύσος}}）；但若源自阿提卡的米乌斯——波洛多罗斯说他是在狩猎中被杀——也无妨；我不嫉妒你们的奥秘享有葬礼般的荣耀。你们也可以换一种方式理解mysteria，即mytheria（狩猎故事），交换两个词的字音；因为这类故事确实追猎最野蛮的色雷斯人、最愚钝的弗吕吉亚人，以及希腊人中迷信者。\par

那么，愿那在人类中首创此欺骗的人归于灭亡，无论是教导众神之母奥秘的达耳达诺斯，还是创立萨莫色雷斯人狂欢和奥秘的厄提昂，或是那弗吕吉亚的弥达斯——他从奥德律索斯学得狡诈的欺骗，又传授给他的人民。因为我绝不相信那个塞浦路斯的基尼拉斯，他竟敢于将阿佛洛狄忒的淫荡狂欢从黑夜带到光天化日之下，热切地要把本国的一个娼妓神化。另一些人说，阿米塔昂的儿子墨兰普斯将色雷斯的节庆从埃及带入希腊，以歌谣纪念她的哀伤。\par

我要把这些事例作为邪恶的首创者、不虔寓言和致命迷信的父母，他们向人类生活中播下了那邪恶与毁灭的种子——奥秘。\par

如今，时机已到，我要证明他们的狂欢充满欺骗和江湖伎俩。你若已受入教礼，更会嘲笑这些曾被尊崇的寓言。我毫无保留地公布那些隐藏的秘密，不羞于说出你们不羞于崇拜的事。\par

于是有那泡沫所生、塞浦路斯所生、基尼拉斯的宠儿——我指的是阿佛洛狄忒，她喜爱男性生殖器，因为从它们而生，即从被割下的乌拉诺斯的那些器官——那些好色的肢体，被割下后对海浪施暴。如此猥亵的肢体生出相称的果实——阿佛洛狄忒。在庆祝这海上欢愉的仪式中，作为她出生的象征，一团盐和\OriginalTerm{阳具}{phallus}被递给那些入不洁奥秘者。那些入教者向她还钱，如同妓女的嫖客向她付钱。\par

还有得默忒耳的奥秘，宙斯对他母亲的无耻拥抱，以及得默忒耳的愤怒——我不知道今后该称她为母亲还是妻子，为此她被称为布里摩，如所说那样；还有宙斯的恳求、苦胆之饮、祭品心脏的摘取，以及我们不敢命名的事。弗吕吉亚人为尊崇阿提斯、库柏勒和科律班忒斯举行这样的仪式。故事说，宙斯扯下一只公羊的睾丸，拿出来扔在得默忒耳的胸前，从而为他强暴的拥抱付出欺诈的代价，假装割下自己的睾丸。这些仪式的入教象征，若在你们空闲时展示，我知道会使你们发笑，尽管由于揭露而毫无笑意。\Quote{我吃自鼓中，我饮自铙钹中，我携带了克诺斯，我溜进了卧室。}这些记号岂非耻辱？这些奥秘岂非荒谬？\par

若我再补充其余部分呢？\Person{得墨忒耳}成为母亲，\Person{科瑞}被抚育成人。随着时间推移，她的生父——同一位\Person{宙斯}——竟与自己的女儿\Person{斐瑞法塔}交合，在母亲\Person{刻瑞斯}之后——忘记了他先前可憎的恶行。\Person{宙斯}既是\Person{科瑞}的父亲，又是其诱奸者，并以龙形无耻地追求她；然而，他的身份被发现了。\Person{萨巴齐奥斯}秘仪中向受印者展示的标记是\Quote{神祇滑过胸膛}——这神祇便是这条爬过受印者胸膛的蛇。这无疑是\Person{宙斯}放纵情欲的明证。\Person{斐瑞法塔}生了一个孩子，尽管是以公牛的形式，正如一位偶像崇拜诗人所说：\par

\Quote{公牛\newline{}龙的父亲，以及公牛的父亲是龙，\newline{}在山丘上牧人隐藏的赶牛刺棒，}\par

你是否希望我讲述\Person{珀耳塞法塔}采花、她的篮子、她被\Person{普路同}（\Person{埃多纽斯}）掳走、大地裂口，以及\Person{欧布勒斯}的猪连同两位女神一起被吞没的故事？为此，在\OriginalTerm{塞斯摩弗里亚节}{Thesmophoria}上，他们操着墨伽拉方言，将猪推出。这则神话故事由妇女们在各城邦以不同的方式庆祝，节日称为\OriginalTerm{塞斯摩弗里亚节}{Thesmophoria}和\OriginalTerm{斯基罗弗里亚节}{Scirophoria}；以多种形式戏剧化\Person{斐瑞法塔}或\Person{珀耳塞法塔}（\Person{普洛塞庇娜}）被掳的事件。\par

\Person{狄俄倪索斯}的秘仪完全是非人性的；因为他尚是孩童时，\Person{枯瑞忒斯}围绕他的摇篮舞动，碰撞兵器，而\Person{提坦}们悄悄潜入，以孩童玩具诱骗他，这些\Person{提坦}在他还是孩童时将他肢解，正如这秘仪的吟游诗人、色雷斯的\Person{俄耳甫斯}所说：\par

\Quote{椎体、陀螺、摇动四肢的响铃，\newline{}以及来自清音\Person{赫斯珀里得斯}的美丽金苹果。}\par

而展示这些神秘仪式的无用象征物，对于谴责并非无用。它们是骰子、球、环、苹果、陀螺、镜子、一簇羊毛。\par

继续我们的叙述，\Person{雅典娜}（\Person{密涅瓦}）取出了\Person{狄俄倪索斯}的心脏，因而被称为\Person{帕拉斯}，源于心脏的颤动；而那些将他肢解的\Person{提坦}们，将鼎置于三脚架上，把\Person{狄俄倪索斯}的肢体投入其中，先煮沸，然后串在签子上，\Quote{置于火上。}但\Person{宙斯}显现，因他是神，迅速察觉到了正在烹煮的肉块的气味——那种你们的神祇一致同意作为他们应得份额的气味——用霹雳袭击\Person{提坦}，并将\Person{狄俄倪索斯}的肢体交给他的儿子\Person{阿波罗}去埋葬。而他——因未违抗\Person{宙斯}——将肢解的尸身带到\Place{帕耳那索斯山}，安葬在那里。\par

你若想考察\Person{科律班忒斯}的狂欢，那么要知道，他们杀死了自己的第三个兄弟，用紫布覆盖死者的头，为其加冠，以矛尖挑起，埋葬在\Place{奥林波斯山}脚下。这些秘仪简言之就是谋杀与葬礼。而主持这些仪式的司祭，被那些负责命名的人称为圣仪之王，给这悲剧事件增添了额外的怪异，禁止将有根的欧芹放在桌上，因为他们认为欧芹是从流淌的\Person{科律班忒斯}之血中长出的；正如妇女们在庆祝\OriginalTerm{塞斯摩弗里亚节}{Thesmophoria}时，禁食落地的石榴籽，因为认为石榴是从\Person{狄俄倪索斯}血滴中长出的。这些\Person{科律班忒斯}也被称为\OriginalTerm{卡比里}{Cabiric}；而仪式本身被宣布为\OriginalTerm{卡比里秘仪}{Cabiric mystery}。\par

因为那两个一模一样的杀兄弟者，取出了存放\Person{巴克斯}阳具的盒子，将其带到埃特鲁里亚——真正贩卖体面货物的人。他们流亡在那里，致力于传授他们迷信的宝贵教导，并将阳具象征物和盒子呈给\Place{第勒尼安人}崇拜。有人不无可能地认为，\Person{狄俄倪索斯}因此被称为\Person{阿提斯}，因为他被阉割了。而\Place{第勒尼安人}是野蛮人，竟被引入如此污秽的仪式，这有什么可惊奇的？因为在\Place{雅典}人中间，在整个希腊——我羞于启齿——关于\Person{得墨忒耳}的可耻传说仍然根深蒂固！因为\Person{得墨忒耳}在寻找女儿\Person{科瑞}时，在\Place{阿提卡}的一个地方\Place{厄琉息斯}附近疲惫不堪，悲痛地坐在一口井边。这对于受印者至今仍是禁止的，以免他们似乎模仿哭泣的女神。当时土著居民占据了\Place{厄琉息斯}：他们的名字是\Person{鲍玻}、\Person{杜扫勒斯}、\Person{特里普托勒摩斯}；此外还有\Person{欧摩尔波斯}和\Person{欧布勒斯}。\Person{特里普托勒摩斯}是牧牛人，\Person{欧摩尔波斯}是牧羊人，\Person{欧布勒斯}是牧猪人；从他们产生出\Person{欧摩尔波斯族}和\Person{传令官族}——一个\Person{圣祭司}的族系——在\Place{雅典}兴盛。\par

那么（因为我将不回避叙述），\Person{鲍玻}殷勤地接待了\Person{得墨忒耳}，递给她一杯提神饮品；女神拒绝饮用，毫无喝的兴趣（因为她非常悲伤），\Person{鲍玻}恼怒，认为自己被轻视，便解开下体，向女神展示其裸体。\Person{得墨忒耳}见了甚为欢喜，勉强接过了饮品——我再说一遍，因这景象而喜悦。这些就是\Place{雅典}人的秘密秘仪；\Person{俄耳甫斯}记录了这些。我将引用\Person{俄耳甫斯}的原话，使你们拥有这秘仪的伟大权威本人作为这桩丑事的证据：\par

\Quote{她如此说着，拉下自己的衣裳，\newline{}露出那不宜提及的身体形态，\newline{}\Person{鲍玻}亲手在乳房下脱去衣物。\newline{}女神于是温和地笑了，心中暗笑，\newline{}接过那闪光的杯，其中盛着饮品。}\par

以下是厄琉息斯秘仪的标记：\Emph{我已禁食，我已饮杯；我已从盒中领取；完成后，我将其放入篮子，再从篮子放入箱子。}真是美妙的景象，与女神相称；适于黑夜、火焰以及厄瑞克透斯族人（\Person{Erechthidæ}）和其他希腊人的秘仪——他们自大而愚蠢，\Quote{死后有他们料想不到的命运等着他们。}确实，厄弗所的赫辣克利托（\Person{Heraclitus}）针对这些人预言道，称他们为\Quote{夜游者、术士、酒神狂女、勒奈狂欢者、受秘仪者。}他威胁他们将面临死后之事，预言他们将受火刑。因为在人眼中被视为秘仪的东西，他们却亵渎地庆祝。那么，律法和意见都是虚无的。龙的秘仪是一个骗局，它表面上虔诚地庆祝根本不是秘仪的秘仪，并以虚假的虔诚遵守亵渎的仪式。这些神秘的箱子是什么？——因为我必须揭露其神圣之物，说出不宜言说之事。难道不是芝麻饼、金字塔形饼、球形和扁平饼（布满浮雕）、盐块、以及一条蛇——\OriginalTerm{巴萨柔斯·狄俄尼索斯}{Dionysus Bassareus}的象征吗？此外，还有石榴、枝条、杖、常春藤叶？此外，还有圆饼和罂粟籽？还有，有忒弥斯（\Person{Themis}）不可言说的象征：墨角兰、灯、剑、女人的梳子——这是对\Emph{女性器官}（\OriginalTerm{muliebria}{muliebria}）的委婉及神秘表达。\par

哦，无耻的脸红！曾几何时，黑夜是沉默的，是节制者欢乐的帷幕；但现在，对于受秘仪者，圣夜却成了放荡仪式的告密者；火炬的闪耀揭露了邪恶的放纵。熄灭火焰，哦，圣师（\Person{Hierophant}）；尊敬火炬，哦，火炬手（\Person{Torch-bearer}）。那光暴露了伊阿科斯（\Person{Iacchus}）；让你的秘仪受到尊崇，命令狂欢在黑夜与黑暗中隐藏。\par

火不伪装；它暴露并惩罚它所受命之事。\par

这就是无神论者的秘仪。我有理由称那些不认识真天主、无耻崇拜一个被泰坦撕碎的男孩、一个受苦的女人、以及因羞耻而实际不可言说的身体部位的人为无神论者；他们被双重不敬所束缚：首先，他们不认识天主，不承认那真实存在的一位为天主；其次，他们错误地认为不存在的存在，称那些没有真实存在、或更恰当地说根本不存在、只有名字的东西为神。因此，宗徒责备我们说：\BibleQuote{你们在应许的盟约上是局外人，没有希望，在世上没有天主。}\BibleRef{弗 2:12}\par

完全归功于那位西徐亚人的国王，无论阿那卡尔西斯（\Person{Anacharsis}）是谁，他射箭处死了他的一名臣民，此人因效仿库梓科斯（\Place{Cyzicus}）居民所行的众神之母的秘仪而在西徐亚人中敲鼓、摇钹（钹挂在脖子上，如同库柏勒（\Person{Cybele}）的司铎），谴责他在希腊人中变得女性化，并成了向其他西徐亚人传授女性化疾病的师傅。\par

为此（我绝不能隐瞒），我不禁感到惊奇，何以阿格里真托的欧赫美鲁斯（\Person{Euhemerus}）、塞浦路斯的尼卡诺尔（\Person{Nicanor}）、狄亚戈拉斯（\Person{Diagoras}）、墨洛斯的希波（\Person{Hippo}）、以及那位名叫特奥多罗斯（\Person{Theodorus}）的库勒涅人等许多人，他们过着清醒的生活，比世人更清晰地看到当时关于那些神祇的流行错误，竟被称为无神论者；因为即使他们没有达到真理的知识，至少也怀疑了普遍观点的错误；这种怀疑不是微不足道的种子，而成为真正智慧的萌芽。其中一人如此指责埃及人：\Quote{如果你们相信他们是神，就不要哀悼或哭号他们；如果你们哀悼并哭号他们，就不要再把他们当作神。}另一人拿起一个赫剌克勒斯（\Person{Hercules}）的木像（他很可能正在家里煮饭），说道：\Quote{来吧，赫剌克勒斯；现在是时候为我们完成这第十三件工作，正如你为欧律斯透斯（\Person{Eurystheus}）完成了十二件，为狄亚戈拉斯（\Person{Diagoras}）预备好这个。}然后把它像木头一样扔进火里。因为无知的两极是无神论和迷信，我们必须努力避免。你们难道没有看见真理的圣师梅瑟（\Person{Moses}）命令，不可让阉人、被阉割者或娼妓之子进入会众吗？前两者暗示一种不敬的习俗，使人被剥夺了神圣的力量和男性气概；第三者则指那些以许多假神取代唯一真天主的人——如同娼妓之子，因不知真正的父亲，可能会声称许多名义上的父亲。\par

人与上天之间有一种与生俱来的原始共融，但因无知而蒙蔽，如今终于瞬间从黑暗中跃出，光芒四射；正如有人以下列诗句所表达的：——\par

\Quote{看这崇高的、无限的以太，\newline{}以它湿润的双臂环绕大地。}\par

以及在这些诗句中：——\par

\Quote{哦，祢以大地为战车，并在大地中安坐，\newline{}无论祢是谁，使我们的努力无法目睹祢。}\par

以及诗人们的子孙所歌咏的其他一切。\par

然而，错误的、偏离正道的、且肯定有害的思想，使具有天上来源的人类转离了天上的生活，将他们按在地上，诱使他们依附于地上的事物。因为有些人被对天空的沉思所迷惑，只信赖自己的视觉，当他们观看星辰的运行时，立刻被惊叹所攫取，将其神化，因星辰的运动而称它们为神（\OriginalTerm{神}{\GreekText{θεός}}来自\OriginalTerm{奔跑}{\GreekText{θεῖν}}）；并崇拜太阳——例如印度人；崇拜月亮——如弗里吉亚人。另一些人，采摘大地所生植物的有益果实，称谷物为得墨忒耳（\Person{Demeter}），如雅典人；称葡萄藤为狄俄尼索斯（\Person{Dionysus}），如忒拜人。还有一些人，考虑到邪恶的惩罚，将其神化，崇拜各种形式的报应和灾难。因此，厄里倪厄斯（\Person{Erinnyes}）、欧墨尼德斯（\Person{Eumenides}）、赎罪神祇以及犯罪的审判者和复仇者，都是悲剧诗人的创造。\par

甚至在诗人之后，你们的一些哲学家也把你们胸中的情感形态，如恐惧、爱情、喜悦、希望，塑造成偶像；例如古代的\Person{厄庇美尼德}，他在雅典筑起了侮辱和无耻的祭坛。人类神化的其他对象则起源于事件，并被赋予形体，如\OriginalTerm{狄刻}{Dike}、\OriginalTerm{克洛托}{Clotho}、\OriginalTerm{拉刻西斯}{Lachesis}、\OriginalTerm{阿特洛波斯}{Atropos}、\OriginalTerm{赫玛墨涅}{Heimarmene}、\OriginalTerm{奥克索}{Auxo}和\OriginalTerm{塔洛}{Thallo}，这些都是阿提卡的女神。还有第六种引入错误和制造神祇的方式，即他们列举十二位神祇，这些神的出生是\Person{赫西俄德}在\WorkTitle{神谱}中歌唱的主题，也是\Person{荷马}在有关神祇的一切叙述中提到的。最后一种方式（总共有七种）起源于天主对人类的神恩。因为人们不明白是天主在善待我们，他们便发明了救主，如\OriginalTerm{狄奥斯库若兄弟}{Dioscuri}、避祸者\OriginalTerm{赫拉克勒斯}{Hercules}和治疗者\OriginalTerm{阿斯克勒庇俄斯}{Asclepius}。这些是从真理中滑落的有害偏差，将人从天上拉下，投入深渊。我想彻底显明你们这些神祇的真相，好使你们如今终于能放弃幻觉，飞速逃回天上。\BibleQuote{因为我们从前也都在过犯中，和其余的人一样，本是义怒之子；然而富于慈悲的天主，由于祂爱我们的大爱，当我们死在过犯中时，竟使我们同基督一起活过来。} \BibleRef{弗 2:3} 因为圣言是活的，既与基督同葬，也同天主一起被高举。但那些仍未相信的人被称为义怒之子，为义怒所养育。我们这些从错误中被拯救出来、回归真理的人，不再是义怒的养育者。因此，我们这些从前是无法无天之子的人，如今通过圣言的慈爱，成了天主的子女。\par

但你们的一位诗人，\Place{阿格里真托}的\Person{恩培多克勒}，前来说道：\par

\Quote{因此，因严重的灾祸而心烦意乱，\newline{}你永不会舒缓你灵魂的悲惨痛苦。}\par

关于你们神祇的大部分传说都是虚构和捏造的；而那些被认为已经发生的事件，则记载了过着放荡生活的卑劣之人：\par

\Quote{你们行走在骄傲和疯狂中，\newline{}离开了正直和笔直的道路，你们走入了\newline{}荆棘和蒺藜。为何你们流浪？\newline{}愚昧的人，停止与凡人交往；\newline{}离开黑暗之夜，把握光明。}\par

这些劝告是\Person{西比拉}——她既是先知又是诗人——吩咐我们的；真理也吩咐我们这样做，它揭去那群神祇可怕而带有威胁的面具，通过显示名称的相似性，驳斥了关于他们的轻率意见。因为有些人算出了三位\OriginalTerm{朱庇特}{Jupiter}：一位是\Place{阿尔卡狄亚}的以太之子，另两位是\OriginalTerm{克罗诺斯}{Kronos}的儿子；其中一位在\Place{克里特}，另外在\Place{阿尔卡狄亚}。还有人算出了五位\OriginalTerm{雅典娜}{Athene}：第一位是\Place{雅典}的\Person{赫淮斯托斯}之女；第二位是\Place{埃及}的\Place{尼罗河}之女；第三位是战争的发明者，\Person{克罗诺斯}之女；第四位是\Person{宙斯}之女，\Place{墨塞尼亚}人称她为\OriginalTerm{科吕法西亚}{Coryphasia}，以其母亲的名字；此外还有\Person{帕拉斯}和\Person{提坦尼斯}之女，\Place{俄刻阿诺斯}之女，她邪恶地杀了自己的父亲，并像穿羊皮一样穿上她父亲的皮。此外，\Person{亚里士多德}称第一位\OriginalTerm{阿波罗}{Apollo}为\Person{赫淮斯托斯}和\Person{雅典娜}之子（因此雅典娜不再是处女）；第二位是在\Place{克里特}，\Person{科里巴斯}之子；第三位是\Person{宙斯}之子；第四位是\Place{阿尔卡狄亚}的\Person{西勒诺斯}之子（阿尔卡狄亚人称他为\OriginalTerm{诺弥乌斯}{Nomius}）；此外，他还列举了\Place{利比亚}的\Person{阿波罗}，\Person{阿蒙}之子；文法家\Person{狄迪摩斯}又加了第六位，\Person{玛格涅斯}之子。如今有多少位\Person{阿波罗}呢？他们数不胜数，都是凡人，与前面提到的一样，是所有同类的助人者。我又何必提那些多位\Person{阿斯克勒庇俄斯}、所有被算出来的\OriginalTerm{墨丘利}{Mercury}，或传说中\OriginalTerm{伏尔甘}{Vulcan}呢？我用这些众多名称淹没你们的耳朵，岂不显得过分吗？\par

无论如何，你们神祇的故乡，以及他们的技艺、生活，尤其还有他们的坟墓，都证明他们是人。例如，战神\OriginalTerm{马尔斯}{Mars}——诗人们对他推崇备至：\par

\Quote{马尔斯，马尔斯，人类的祸根，染血的城墙摧毁者，}\par

——这位神祇，总是变节且不可和解，正如\Person{厄庇卡摩斯}所说，是一个\Place{斯巴达}人；\Person{索福克勒斯}知道他是\Place{色雷斯}人；其他人说他是一个\Place{阿尔卡狄亚}人。\Person{荷马}说，这位神祇被捆绑了十三个月：\par

\Quote{马尔斯曾受苦；被\Person{阿洛欧斯}的儿子们，\newline{}\Person{俄托斯}和\Person{厄菲阿尔忒斯}，强力捆绑，\newline{}他在铜锁链中躺了十三个月。}\par

祝\Place{卡里亚}人好运，他们向他献狗为祭！愿\Place{西徐亚}人永不停止献驴为祭，正如\Person{阿波罗多洛斯}和\Person{卡利马科斯}所记载：\par

\Quote{福玻斯向\Place{许珀耳玻瑞亚}人显现恩惠，\newline{}然后他们向他献上驴为祭。}\par

同一作者在另一处写道：\par

\Quote{肥美的驴肉祭品使福玻斯喜悦。}\par

\Person{赫淮斯托斯}，被\Person{朱庇特}从\Place{奥林帕斯}山的神圣门槛扔下，落在\Place{利姆诺斯}岛上，学会了铜匠手艺，双脚残疾：\par

\Quote{他摇晃的膝盖在重压下弯曲。}\par

你们神中还有一位医生，而不仅仅是铜匠。那医生贪爱黄金；他的名字是\Person{阿斯克勒庇俄斯}。我请你们自己的诗人、\Place{波奥提亚}的\Person{品达}作证：\par

\Quote{连他也被手中闪闪发光的金子，\newline{}相当于一笔可观的酬金，说服了\newline{}从死神手中救回一个已为死所掳的人；\newline{}但\Person{克罗诺斯}之子\Person{朱庇特}，用霹雳击中两者，\newline{}迅速从他们胸膛夺走呼吸，\newline{}他的闪电霹雳封住了他们的命运。}\par

还有\Person{欧里庇得斯}：\par

\Quote{因为\Person{宙斯}对我儿子\Person{阿斯克勒庇俄斯}的死有罪，\newline{}他将闪电火焰投向他的胸膛。}\par

因此，他被闪电击中，躺在\Place{库诺苏里斯}地区。\Person{菲洛科罗斯}也说，在\Place{特诺斯}岛，\Person{波塞冬}被当作医生崇拜；\Person{克罗诺斯}定居在\Place{西西里}，并埋葬在那里。\Place{图里亚}的\Person{帕特罗克洛斯}和\Person{小索福克勒斯}在三部悲剧中讲述了\Person{狄奥斯库若兄弟}的故事；而这两位\Person{狄奥斯库若}不过是两个凡人，如果\Person{荷马}值得信赖的话：\par

\Quote{……但他们躺在丰饶的土地之下，\newline{}在拉栖代梦，他们的故土。}\par

此外，撰写塞浦路斯诗篇的人说卡斯托尔是有死的，命运已为他定下死亡；而波鲁克斯是不死的，是玛尔斯的后裔。这是他以诗人的方式虚构的。但荷马更可信，他如上所述论及两位狄奥斯库里，并且还证明赫拉克勒斯只是一个幻影：——\par

\Quote{赫拉克勒斯，精通大能事迹的人。}\par

因此，赫拉克勒斯在荷马本人看来只是一个有死之人。哲学家希罗尼穆斯描述了他的体格：高大、毛发竖立、强壮；狄凯阿科斯说他身材方正、肌肉发达、肤色黝黑、鹰钩鼻、灰眼睛、长发。这位赫拉克勒斯活了五十二年后，在奥塔山的火葬堆上被焚烧，结束了生命。\par

至于缪斯，阿尔坎德尔称她们是宙斯和谟涅摩叙涅的女儿，其他诗人和作家也神化并崇拜她们——那些缪斯，为了纪念她们，许多城邦已经建立了博物馆，她们原是侍女，被马卡尔的女儿墨伽克洛雇佣。这位马卡尔统治莱斯博斯岛，经常与妻子争吵；墨伽克洛为母亲感到烦恼。她为母亲会做什么呢？于是她雇佣了那些侍女，数目众多，按照伊奥利亚人的方言称她们为缪塞。她教导她们歌唱古时的事迹，并动听地弹奏里拉琴。她们通过不断弹奏里拉琴和甜美歌唱，安抚了马卡尔，制止了他的坏脾气。因此，墨伽克洛出于对她们的感激，为母亲竖立了铜柱，并命令在所有神庙中尊崇她们。这就是缪斯的来历。此记载出自莱斯博斯岛的米尔西卢斯。\par

现在，请听你们诸神的爱情，以及他们放纵情欲的难以置信的故事，还有他们的伤痕、捆绑、欢笑、争斗、奴役和宴饮；此外，还有他们的拥抱、眼泪、苦难和淫秽的欢愉。请呼波塞冬，以及被他玷污的一群少女：安菲特里忒、阿弥摩涅、阿洛佩、墨拉尼佩、阿尔库俄涅、希波托厄、基俄涅，以及成千上万的其他女子；尽管有这么多，你们的波塞冬的情欲仍未满足。\par

请呼阿波罗；这是福玻斯，既是圣洁的先知，又是良善的顾问。但斯忒罗佩不会这么说，埃托乌萨、阿尔西诺厄、宙克西佩、普罗托厄、玛耳皮萨、许普西皮勒也不会。因为只有达佛涅逃脱了这位先知和诱惑。\par

最重要的是，按照你们的说法，让诸神和人类的父亲亲自来吧，他是如此沉溺于肉欲，以致贪恋一切，并在一切之上满足他的欲望，就像特姆伊特人的山羊一样。哦，荷马，你的诗篇令我赞叹！\par

\Quote{他说了，并以他阴影般的眉毛点头；\newline{}不朽头颅上的神妙发卷拂动，\newline{}整个奥林波斯山在他的点头下颤抖。}\par

哦，荷马，你使宙斯显得可敬；你赋予他的点头是最庄严的。但只给他看一条女人的腰带，宙斯就被暴露了，他的头发被玷污了。你们的宙斯堕落到了何等程度，竟与阿尔克墨涅度过了如此多的纵欲之夜？因为这九个夜晚对这个贪得无厌的怪物来说都不算长。相反，整个一生都不足以满足他的欲望；这样他才能为我们生出那避祸的神祇。\par

赫拉克勒斯，宙斯之子——一个真正的宙斯之子——是那个长夜的产物，他历经艰辛在很长时间内完成了十二项劳动，但在一个夜里玷污了忒斯提俄斯的五十个女儿，因而同时是这么多处女的淫乱者和新郎。因此，诗人们称他为残忍的恶棍和邪恶的流氓，这不无道理。叙述他的各种通奸和引诱男孩的行为将不胜枚举。因为你们的诸神甚至没有远离男童，一个爱上了许拉斯，另一个爱上了许阿铿托斯，另一个爱上了珀罗普斯，另一个爱上了克律西波斯，另一个爱上了伽倪墨得斯。让这样的神祇被你们的妻子崇拜吧，让她们祈求她们的丈夫也像这样——如此节制；这样，效仿同样的行为，他们就能像神一样。让你们的男孩受训练去崇拜这样的神，以便他们长大后，身上带着从神那里接受的受诅咒的淫乱形象，成为男人。\par

但也许只有男性神祇在性放纵上如此冲动。\par

\Quote{女神们因羞耻而各自待在屋里，}荷马说；女神们因贞洁而脸红，不敢看犯了通奸的阿芙洛狄忒。但她们更加放荡不羁，被通奸的锁链束缚：厄俄斯因提托诺斯而蒙羞，塞勒涅因恩底弥翁而蒙羞，涅瑞伊斯因埃阿科斯而蒙羞，忒提斯因珀琉斯而蒙羞，得墨忒耳因伊阿宋而蒙羞，珀耳塞福涅因阿多尼斯而蒙羞。阿芙洛狄忒因阿瑞斯而蒙羞后，又去找喀倪拉斯，嫁给安喀塞斯，为法厄同设下陷阱，并爱上阿多尼斯。她与牛眼的朱诺竞争；女神们为了苹果而解衣，赤身裸体地出现在牧羊人面前，让他判决谁最美丽。\par

但来吧，让我们简要地巡视各项竞赛，并废除那些在坟墓举行的庄严集会：地峡运动会、涅墨亚运动会、皮提亚运动会，最后是奥林匹亚运动会。在皮托，皮提亚巨蛇被崇拜，而蛇的节日集会被称为皮提亚。在地峡，大海吐出一块可怜的垃圾；地峡运动会哀悼墨利克尔塔。\par

在涅墨亚，另一个——一个小男孩，阿尔克摩洛斯——被埋葬；这孩子的殡葬竞技会称为涅墨亚竞技会。皮萨是弗里吉亚战车御者的坟墓，哦，各族裔的希腊人；而奥林匹亚竞技会，不过是珀罗普斯的殡葬祭献，非狄亚斯的宙斯将其据为己有。神秘仪式，很可能，也是为死者举行的竞技会；神谕亦是如此，两者后来都成为公开的。但萨格拉和阿提卡的阿里穆斯的秘仪仅限于雅典。而那些献给狄俄尼索斯的竞技和\Emph{阳具像}，则是世间的耻辱，以其致命的影响渗透生命。因为狄俄尼索斯热切渴望下到冥府，却不知路径；一个名叫普洛绪姆诺斯的人提议告诉他，但非无报酬。那报酬是可耻的，尽管在狄俄尼索斯看来并非如此：他向狄俄尼索斯求取一种阿佛洛狄忒式的恩惠作为报酬。神并不勉强答应他的请求，并承诺若他返回便履行，且以誓言确认。得知路径后，他离去又返回：他没有找到普洛绪姆诺斯，因为他已死了。为了履行对这位情人的承诺，他冲向他的坟墓，燃起反常的情欲。他砍下一根随手可得的无花果树枝，塑造了阳具，从而向死者履行了诺言。作为这一事件的秘仪纪念，\Emph{阳具像}被高举在各城市中以尊崇狄俄尼索斯。\Quote{因为若他们不举行敬拜狄俄尼索斯的游行，并歌唱最无耻的颂扬生殖器的歌曲，一切都会出错，}赫拉克利特说。这就是那位冥王和狄俄尼索斯，因他们而陷入狂热、扮演酒神信徒——依我之见，与其说是为了醉酒，不如说是为了无耻的仪式。因此，顺理成章，凡成为自己情欲奴隶者，便是你们的神！\par

此外，如同拉刻代蒙人中的希洛人，阿波罗在斐赖受制于阿德墨托斯的奴役，赫拉克勒斯在萨迪斯受制于翁法勒。波塞冬为拉俄墨冬服役；阿波罗也是如此，他像一个无用的仆人，无法从先前的主人那里获得自由；那时特洛伊城墙就是由他们为弗里吉亚人建造的。荷马也不羞于提及雅典娜手持金灯出现在奥德修斯面前。我们还读到阿佛洛狄忒，像一个轻浮的使女，为海伦取来并安置座位，让她坐在那奸夫对面，以便引诱他。\par

帕尼亚西斯也告诉我们，除了那些充当仆人的神之外，还有许多别的神，这样写道：——\par

\Quote{得墨忒耳曾受奴役，著名的瘸腿神也是如此；\newline{}波塞冬也曾如此，还有银弓的阿波罗，\newline{}与一个凡人同住一年。凶猛的阿瑞斯\newline{}也曾如此，在他父亲的强迫之下。}\par

如此等等。\par

与此相应，我还要向你们提出那些多情而肉欲的神祇，他们在各方面都有人类的情感。\par

\Quote{因为他们有凡人的身体。}\par

荷马最清楚地表明了这一点，他描述阿佛洛狄忒因受伤而发出高声尖叫；并描述最善战的阿瑞斯本人被狄俄墨得斯刺伤腹部。波勒摩也说雅典娜被俄耳尼托斯所伤；不仅如此，荷马说冥王普路托也被赫拉克勒斯的箭射中；帕尼亚西斯叙述太阳神赫利俄斯的光线被赫拉克勒斯的箭射中；同一位帕尼亚西斯叙述，同一位赫拉克勒斯使婚姻女神赫拉在多沙的皮洛斯受伤。索西比乌斯也叙述赫拉克勒斯被希波科翁的儿子们伤到手。若有伤口，便有血。因为诗人的\OriginalTerm{神血}{ichor}比血更令人厌恶；血的腐败称为神血。因此必须提供他们所需的治疗和食物。于是提到餐桌、饮酒、欢笑和交媾；因为人若不死，无所需求，永不衰老，就不会投身于爱情、生育子女或睡眠。朱庇特本人，当他在阿卡狄亚人吕卡翁家做客时，在埃塞俄比亚人中分享了人肉餐——一餐极不人道且被禁止的。因为他无意中以人肉果腹；这位神不知道，他的款待者阿卡狄亚人吕卡翁已杀了自己的儿子（名叫倪克提摩斯），烤熟后端到宙斯面前。\par

这就是朱庇特，善的、预言的、好客的、保护恳求者的、慈祥的、预兆的作者、伸冤者；而非不义的、践踏正义与法律的、不敬的、不人道的、暴力的、诱惑者、奸淫者、好色者。但或许当他如此时，他本是人；而现在这些神话在我们看来似乎已经古老了。宙斯不再是蛇、天鹅、鹰，也不是淫荡的人；神不再飞翔，不爱少年，不亲吻，不施暴，尽管仍有许多美丽的女子，比勒达更美，比塞墨勒更鲜嫩，还有容貌和品行比弗里吉亚牧童更好的少年。那鹰如今在何处？那天鹅如今在何处？宙斯本人又在何处？他已随着他的羽毛老去；因为他至今仍不悔改他的风流韵事，也未学会节制。神话已揭露在你们面前：勒达死了，天鹅死了。寻找你们的朱庇特。不要搜索天堂，而要搜索大地。克里特人，在他的国土上他被埋葬，会将指给你们看——我指的是卡利马科斯在他的颂歌中：——\par

\Quote{因为你的坟墓，哦，王啊，\newline{}克里特人建造了！}\par

因为宙斯死了，不要悲伤，正如勒达死了，天鹅、鹰、浪荡子和蛇都死了。如今即使迷信的人，似乎虽不情愿，却真实地开始明白他们对诸神的错误看法。\par

\Quote{因为你们并非出自古老的橡树，或岩石，\newline{}而是出自人类。}\par

但此后不久，他们将被发现不过只是橡树和石头。斯塔菲罗斯说，一个阿伽门农在斯巴达被当作朱庇特崇拜；法诺克勒斯在其 \WorkTitle{勇敢者与美者} 一书中述说，赫伦人的国王阿伽门农为纪念其友阿尔根努斯，修建了阿尔根尼亚的阿佛洛狄忒的庙宇。一位被称为\Quote{被勒死}的阿尔忒弥斯受到阿卡狄亚人崇拜，如同卡利马科斯在其 \WorkTitle{原因集} 中所说；而在梅提姆纳，另一位阿尔忒弥斯，即康迪利提斯·阿尔忒弥斯，也获享神荣。还有另一位阿尔忒弥斯的庙宇——波达格拉（即痛风）阿尔忒弥斯——在拉科尼亚，根据索西比奥斯所说。波勒蒙提到一个打哈欠的阿波罗神像；还有另一尊在埃利斯受尊崇的狂饮阿波罗神像。此外，埃利斯人向驱蝇者宙斯献祭；罗马人向驱蝇者赫拉克勒斯、向热病和恐惧（他们也视为赫拉克勒斯的从者）献祭。（我略过阿尔戈斯人，他们崇拜开启坟墓的阿佛洛狄忒。）阿尔戈斯人和斯巴达人尊崇咳嗽者阿尔忒弥斯——称为\OriginalTerm{凯利提斯}{\GreekText{χελύτις}}，从\OriginalTerm{\GreekText{κελύττειν}}{\GreekText{κελύττειν}}一词而来，这词在他们的语言中意为咳嗽。\par

你认为这些细节是从哪里引来的？此处所举的只是你们自己提供的事例；而你们似乎并未认出自己的作家，我呼召他们作证反对你们的不信。可怜虫！你们用不敬的戏谑充斥着你们整个的生命——这生命其实是毫无生命的！\par

难道在阿尔戈斯没有一位秃头宙斯被崇拜？在塞浦路斯另一位复仇者宙斯？难道阿尔戈斯人不向\Quote{护行者}阿佛洛狄忒献祭？雅典人不向\Quote{交际花}阿佛洛狄忒？叙拉古人不向\OriginalTerm{卡利皮戈斯}{\GreekText{Καλλίπυγος}}阿佛洛狄忒献祭？尼坎德尔在某处称她为\OriginalTerm{卡利格洛托斯}{\GreekText{Καλλίγλουτος}}（美臀的）。我在此暂不提及\OriginalTerm{狄俄尼索斯·乔伊罗普萨勒斯}{\GreekText{Διόνυσος} \GreekText{Χοιροψάλης}}。西库昂人尊崇这位神祇，他们把他立为女性器官之神——污秽的庇护者——并虔诚地尊他为放荡的创始者。这些就是他们的神明；这些人也正是在嘲弄神明，或者更确切地说，是在嘲弄并侮辱自己。埃及人在他们的城镇和乡村中向无理性的动物奉献神荣，比起崇拜这些神明的希腊人，他们强了多少？\par

因为，如果它们是野兽，它们便不犯奸淫或好色，也不追求任何违反本性的事。至于这些神祇是何种模样，何须再多说？因他们已充分暴露。此外，我刚才提到的埃及人在崇拜对象上有所分歧。锡耶纳人崇拜\OriginalTerm{锦鳚鱼}{\GreekText{βραχίων}}；埃勒凡廷的居民崇拜\OriginalTerm{迈奥特斯}{\GreekText{μαιώτης}}——这是另一种鱼；俄克叙林库斯人同样崇拜一种以其国家命名的鱼。赫拉克利奥波利人崇拜\OriginalTerm{獴}{\GreekText{ἰχνεύμων}}；赛斯人和底比斯人崇拜羊；琉科波利人崇拜狼；库诺波利人崇拜狗；孟斐斯人崇拜阿匹斯；门德斯人崇拜山羊。而你们，完全胜过埃及人（我宁愿说更差），你们一生没有一天停止嘲笑埃及人，可你们中有些人对待野兽又如何？因为你们中，帖撒利人依照古老习俗向鹳鸟致以神性的敬礼；底比斯人则向鼬鼠致敬，因它们在赫拉克勒斯出生时提供了帮助。另外，不是有报告说帖撒利人崇拜蚂蚁吗？因为他们得知宙斯曾化作蚂蚁与克勒托尔之女欧律墨杜萨交合，生下密尔弥冬。波勒蒙也述说，特洛阿德的居民崇拜当地的鼠类，他们称之为\OriginalTerm{斯敏托伊}{\GreekText{Σμινθοί}}，因为这些老鼠啃断了敌人弓上的弦；据此，阿波罗得到了\Quote{斯敏提乌斯}的称号。赫拉克利得斯在其著作 \WorkTitle{关于阿卡纳尼亚神庙的建造} 中说，在阿克提乌姆海角及阿克提乌姆的阿波罗神庙所在之处，他们向苍蝇献上一头公牛为祭。\par

我也不会忘记萨摩斯人：萨摩斯人，如欧福里翁所说，尊崇羊。我也不会忘记居住在腓尼基的叙利亚人，其中一些人敬拜鸽子，另一些人敬拜鱼，其过度敬拜不下于埃利斯人敬拜宙斯。既然你们所崇拜的不是神，在我看来有必要弄清那些被你们列为第二等级（仅次于诸神）的，是否真是恶魔。因为，如果贪婪与不洁的是恶魔，那么在本土恶神中可以发现成群的神圣荣誉已遍布你们的城市：库特诺斯人中的墨涅得摩斯；忒诺斯人中的卡利斯塔戈拉斯；提洛人中的阿尼乌斯；拉科尼亚人中的阿斯特拉巴库斯；在法勒鲁姆，一位船首的英雄被崇拜；而在波斯战争高潮时，皮提亚女祭司命令普拉提亚人向安德罗克拉忒斯、德谟克拉忒斯、库克勒乌斯和勒乌科献祭。其他众多恶魔可被任何稍加留意的人发掘出来。\par

\Quote{因为在养育万物的大地上，有无数\newline{}不朽的恶魔，是会言说之人的守护者。}\par

这些守护者是谁？波奥提亚人啊，不要吝于道出。岂不显然就是我们提到的那些，以及更著名的那些，伟大的恶魔：阿波罗、阿尔忒弥斯、勒托、得墨忒耳、科瑞、普路同、赫拉克勒斯，乃至宙斯本人？\par

但他们是从逃跑中保护我们，阿斯克拉人啊，或者也许是保护我们远离犯罪，因为他们自己从未尝试过犯罪！若然，那谚语便恰到好处地说：\par

\Quote{未曾受教诫的父亲，却诫管教儿子。}\par

如果这些是我们的守护者，那并非因为他们对我们怀有任何善意，而是专为你们的毁灭，像谄媚者一样，受烟气所诱，吞吃你们的财物。这些恶魔自己也承认他们的贪食，说道：\par

\Quote{因着应得的奠酒和羔羊的脂油，\newline{}我的祭坛一直在他们手中得喂养；\newline{}这样的尊荣始终向我们献上。}\par

如果埃及人的神祇，如猫和鼬鼠，真的获得了说话的能力，它们会说出什么别的话呢？难道不是那宣告它们喜爱香味和烹饪的、荷马式的、诗意的言辞吗？这就是你们的魔鬼、神祇和\OriginalTerm{半神}{demigods}（如果有所谓半神的话，正如有\OriginalTerm{半驴}{demi-asses}（骡子）一样）；因为你们不乏创造亵渎的复合名称的词汇。\par

\SourceRecord{Translated by William Wilson.；Ante-Nicene Fathers；Vol. 2.；Edited by Alexander Roberts, James Donaldson, and A. Cleveland Coxe.；Buffalo, NY: Christian Literature Publishing Co.,；1885.；<http://www.newadvent.org/fathers/020802.htm>.}

% structure: p0001 source=h1 target=section reason=page-title

\section{劝勉外邦人（第三章）}

\Signature{亚历山大的克莱孟}\par

% structure: p0003 source=h2 target=subsection reason=relative-heading-rank; base=h2

\subsection{第三章 祭献诸神的残暴}

那么，现在我们要补充说，你们的诸神是多么不人道、与人类为敌的魔鬼，他们不仅以人的疯狂为乐，而且幸灾乐祸于人的屠杀——时而在竞技场上为争霸而进行的武装竞赛中，时而在无数争荣的战争中，为自己提供快乐的手段，以便能充分满足于对人的谋害。\par

如今，他们像瘟疫一样侵入城邦和民族，要求残忍的祭品。就这样，麦森尼亚人\Person{阿里斯托墨涅斯}为向\Person{伊索梅特的宙斯}献祭而杀了三百人，以为如此数量和质量的全燔祭会有吉兆；其中就有拉西第梦人的国王\Person{塞奥彭波斯}，一个高贵的牺牲。\par

陶里人，即居住在陶里刻松的民族，将任何漂流到他们海岸的陌生人立即祭献给陶里的\Person{阿尔忒弥斯}。\Person{欧里庇得斯}在悲剧舞台上表现了这些献祭。\Person{莫尼穆斯}在其论奇事的著作中记载，在帖撒利的\Place{培拉}，一个阿哈伊亚人被祭献给\Person{佩琉斯}和\Person{喀戎}。\Person{安提克利德斯}在其\WorkTitle{返乡记}中表明，吕克提人（一个克里特种族）曾向\Person{宙斯}献祭人祭；\Person{多西达斯}说，莱斯沃斯人也向\Person{狄俄尼索斯}献上类似的祭品。\Person{皮托克勒斯}在其第三卷\WorkTitle{论和谐}中告诉我们，福西亚人（因为我不放过这样的人）也将一个人作为全燔祭献给陶里亚的\Person{阿尔忒弥斯}。\par

阿提卡的\Person{厄瑞克透斯}和罗马的\Person{马略}祭献了自己的女儿——前者献给\Person{珀耳塞福涅}，正如\Person{德马拉图斯}在其关于悲剧题材的第一卷书中所述；后者献给避邪恶神祇，正如\Person{多罗修斯}在其\WorkTitle{意大利纪事}第一卷中所记。从这些例子来看，魔鬼确实显得仁慈；而敬畏魔鬼的人岂不也相应虔诚？前者被冠以救世主的美名；后者则向那些图谋他们安全的人祈求安全，以为用吉兆向他们祭献，却忘了他们自己在杀人。因为谋杀并不因为在特定地点进行而成为祭献。如果一个人在祭台上或在通往\Person{阿尔忒弥斯}或\Person{宙斯}的大路上杀了人，那并不比因愤怒或贪婪而杀人（其他与前者非常相似的魔鬼）更神圣；这种祭献乃是谋杀和人类的屠宰。那么，哦，人，所有受造物中最有智慧的，为何你们躲避野兽，避开凶猛动物，如果遇到熊或狮子就远远躲开？\par

\Quote{……就如旅人发现，\newline{}盘绕在路中，在山的这边，\newline{}一条致命的蛇，他急忙退缩——\newline{}四肢颤抖，面如土色，}\par

然而，你们既然察觉并理解魔鬼是致命而邪恶的、阴谋家、人类憎恨者和毁灭者，为何不避开它们，或将它们驱逐出去？邪恶者能说出什么真理，或对任何人做什么好事呢？\par

那么，我能轻易证明人比你们这些不过是魔鬼的诸神更优越；例如，我可以指出\Person{居鲁士}和\Person{梭伦}胜过神谕的\Person{阿波罗}。你们的\Person{福玻斯}是贪爱馈赠者，而非爱人者。他背叛了他的朋友\Person{克洛索斯}，忘记了他已得到的报酬（他如此在意自己的名声），领他过\Place{哈吕斯河}走向火刑柱。魔鬼爱人的方式是使人走向（不灭的）火。\par

但哦，人啊，你比\Person{阿波罗}更爱人、更真诚，可怜那绑在柴堆上的人吧。哦\Person{梭伦}，宣告真理；哦\Person{居鲁士}，命令将火熄灭。终于，哦\Person{克洛索斯}，受苦难的教导，变得智慧吧。你所崇拜的是一位忘恩负义者；他接受你的酬报，拿了金子之后却背信弃义。\Quote{再看结局，哦\Person{梭伦}。}告诉你这话的不是魔鬼，而是人。\Person{梭伦}发出的并非模棱两可的神谕。你很容易抓住他。野蛮人啊，当你被放在柴堆上时，你将通过证明发现这个神谕完全是真实的。因此，我不得不诧异，最初迷失的人是被什么动听的道理驱使，向人传布迷信，劝他们崇拜邪恶的魔鬼，无论是\Person{福洛纽斯}还是\Person{墨洛普斯}，或是任何其他为他们修建庙宇和祭坛的人；此外，正如传说，他们是第一个向他们献祭的人。但毫无疑问，在后来的时代，人类为自己发明了崇拜的诸神。毫无疑问，这位据说最古老的神之一的\Person{厄洛斯}，从未被人崇拜，直到\Person{卡耳穆斯}带了一个小男孩，在\Place{阿卡德米}为他立了一座祭坛——这比他满足了的淫欲更为体面；而罪恶的淫乱人们称之为\Person{厄洛斯}之名，这样将无度的淫欲神化。同样，雅典人直到\Person{腓力庇得斯}告诉他们才知道\Person{潘}是谁。\par

那么，正如所料，迷信就这样产生，成为愚妄邪恶的泉源；它后来未被遏制，反而不断增长、如洪水般奔涌，成了众多魔鬼的始作俑者，并且显现在献上馨香祭、设立庄严集会、竖立偶像、建造庙宇——实为坟墓——之中：因为我不愿沉默地放过这些，尽管它们被称为庙宇的庄严名号，实则是坟墓，却得了庙宇之名。但如今你们终当彻底放弃你们的迷信，因以宗教敬意看待坟墓而感到羞愧。在\Place{拉里萨的雅典娜神庙}中，在\Place{卫城}上，有\Person{阿克里西俄斯}的坟墓；在\Place{雅典}的\Place{卫城}上，有\Person{刻克洛普斯}的坟墓，正如\Person{安提奥库斯}在其\WorkTitle{历史}第九卷中所说。\Person{厄里克托尼俄斯}呢？他岂非葬在\Place{波利阿斯神庙}中？\Person{欧摩尔波斯}与\Person{黛拉}之子\Person{伊马洛斯}，岂非葬在\Place{卫城下的厄琉息尼翁圣地}内？\Person{克琉斯}的女儿们，岂非葬在\Place{厄琉息斯}？我何必向你们历数极北地区的妇女？她们名为\Person{许珀罗刻}和\Person{拉俄狄刻}；她们葬在\Place{提洛岛的阿耳忒弥斯庙}中，即在\Place{提洛的阿波罗神庙}内。\Person{勒安德里乌斯}说，\Person{克莱尔科斯}葬在\Place{米利都}的\Place{狄底玛庙}中。遵循\Place{米恩达}的\Person{芝诺}，在此不宜略过\Person{琉科弗吕涅}的坟墓，她葬在\Place{玛格尼西亚的阿耳忒弥斯庙}中；也不宜略过\Place{忒尔墨索斯的阿波罗祭坛}，据说是预言者\Person{忒尔墨修斯}之墓。此外，\Person{阿革萨尔科斯}之子\Person{托勒密}在其关于\Person{菲洛帕托尔}的第一卷书中说，\Person{喀倪剌斯}及其后裔葬在\Place{帕福斯的阿佛洛狄忒庙}中。然而，若要遍历你们视为神圣的坟墓，我的时间就不够用了。如果你们对这些大胆不敬之事毫无羞愧，那你们就是完全死了，正如你们实际所做的那样，信赖死者。\par

\Quote{可怜的人哪，你们遭受的是何等苦难？\newline{}你们的头颅被黑夜的黑暗笼罩。}\par

\SourceRecord{Translated by William Wilson.；Ante-Nicene Fathers；Vol. 2.；Edited by Alexander Roberts, James Donaldson, and A. Cleveland Coxe.；Buffalo, NY: Christian Literature Publishing Co.,；1885.；<http://www.newadvent.org/fathers/020803.htm>.}

% structure: p0001 source=h1 target=section reason=page-title

\section{劝勉外邦人（第四章）}

\Signature{克肋孟．亚历山大}\par

% structure: p0003 source=h2 target=subsection reason=relative-heading-rank; base=h2

\subsection{第四章　偶像崇拜的荒谬与可耻}

如果我再把这些偶像本身拿出来放在你们面前检视，你们在逐一审视时就会发现，你们自幼养成的崇拜人手所造的无意识作品的风俗是多么愚蠢。\par

古时，斯齐提亚人崇拜他们的马刀，阿拉伯人崇拜石头，波斯人崇拜河流。还有一些更古老的部族，在显眼之处竖起木块，立起石柱，这些被称为\OriginalTerm{原始偶像}{Xoana}，源自雕琢其材料。伊卡洛斯的阿耳忒弥斯像无疑是未经雕刻的木头，基泰戎的赫拉像是一棵砍倒的树干；而萨摩斯的赫拉像，据埃特利乌斯说，最初是一块木板，后来在普洛克路斯统治期间才被雕刻成人形。当原始偶像开始被塑成人形时，它们被称为\OriginalTerm{雕像}{Brete}，这词源于\OriginalTerm{人}{Brotos}（人）。在罗马，历史学家瓦罗说，古时玛尔斯的偶像——即被崇拜的神像——是一支矛，当时艺术家尚未涉足这门貌似悦目却有害的艺术；而当艺术繁荣时，谬误也随之增长。你们用石头、木头，简言之，用死物制作了人形的偶像，借此制造了虔敬的赝品，诽谤了真理，如今已是再清楚不过了；但该论点所需的证据仍不可省略。\par

奥林匹亚的宙斯像和雅典的波利亚斯像都是由菲狄亚斯用黄金和象牙制作的，这是众所周知的；奥林匹库斯在《\WorkTitle{萨摩斯纪}》中记述，萨摩斯的赫拉像是由欧几里得的凿子雕成的。因此，不要对所谓的雅典\Quote{可敬之神}有任何怀疑：其中两尊是用名为吕克尼斯的石头制成的，由斯科帕斯雕刻，而卡洛斯雕刻了报告中说置于二者之间的那一尊，正如波勒蒙在致蒂迈欧的第四卷书中所展示的。也不用怀疑吕西亚的帕塔拉的宙斯和阿波罗像，它们出自菲狄亚斯之手，还有与之相伴的狮子；如果如某些人所说，它们是布律克西斯的作品，我也不争辩——你便有了另一位造像者。无论你更喜欢哪一位，都写下吧。此外，在特诺斯被崇拜的九肘尺高的波塞冬和安菲特里忒像，是雅典人特勒西乌斯的作品，据菲洛科鲁斯所述。德米特里乌斯在《\WorkTitle{阿尔戈斯纪}》第二卷中写道，梯林斯的赫拉像的材料是梨木，艺匠是阿尔戈斯。\par

许多人或许会惊讶地得知，被称为\OriginalTerm{从天而降的}{Diopetes}的帕拉狄翁——即从天上掉下来的——据载由狄奥墨得斯和尤利西斯从特洛伊携至得摩丰处，其材料是珀罗普斯的骨头，正如奥林匹亚的朱庇特像用其他骨头——印度野兽的骨头——制成一样。我引述狄奥尼修斯为证，他在《\WorkTitle{周期}》第五部分中记叙了此事。阿佩拉斯在《\WorkTitle{德尔斐纪}》中说，有两个帕拉狄翁，都是人手所造。但为避免有人以为我因无知而略过，我还要补充：雅典的狄俄尼索斯·莫里库斯像是由称为费拉塔的石头制成，是西蒙（欧帕拉摩斯之子）的作品，波勒莫在一封信中如此说。还有两位克里特的雕刻家，我想：名叫斯库勒斯和狄波厄努斯；他们完成了阿尔戈斯的狄奥斯库里像、梯林斯的赫拉克勒斯像，以及西库昂的穆尼基亚的阿耳忒弥斯像。我何必在这些上耽搁，当我能指给您那伟大的神祇本身，并指出他是谁——我们特别清楚地听到他被认为配受敬拜？他们竟敢说他是非人手所造的——我指的是埃及的塞拉皮斯。有些人说，他是西诺普人送给埃及王托勒密·菲拉德尔福斯的礼物，当时西诺普人正遭受饥荒，托勒密从埃及运粮救济他们，因而赢得他们的好感；这偶像原是普路托的像；托勒密接收后，将其置于现在称为拉科提斯的海岬，那里塞拉皮斯神庙备受尊崇，圣所毗邻该地；又说，名妓布利斯蒂基斯死于卡诺普斯，托勒密将她运至此处，葬于上述神庙之下。\par

另有人说，塞拉皮斯是一个本都的偶像，被隆重地运至亚历山大。唯有伊西多尔说，它来自安提约基雅附近的塞琉西亚人，那里也曾发生粮荒，受托勒密接济。但桑顿之子阿特诺多罗斯，试图证明塞拉皮斯是古老的，却不慎滑入错误，证明它是人手所造的像。他说，埃及王塞索斯特里斯征服了大多数希腊民族，返回埃及时带回了许多工匠。于是他下令为他的祖先奥西里斯制作一尊富丽堂皇的雕像；由艺匠布律阿克西斯完成——不是雅典的那位，而是同名的另一位——他在制作中混合了各种材料。他有金、银、铅的碎屑，还有锡；埃及的石头一种不缺，还有蓝宝石、赤铁矿、祖母绿、黄玉的碎片。他将所有这些成分磨碎混合，给混合物染上蓝色，于是雕像呈暗色；然后将整个混合物与奥西里斯和阿皮斯葬礼剩下的颜料混合，塑造了塞拉皮斯，其名称表明它与墓葬有关并由葬礼材料制成，由奥西里斯和阿皮斯结合而成，合称\OriginalTerm{奥西拉皮斯}{Osirapis}。\par

在埃及，又有一位新神祇以隆重的宗教仪式被纳入神班，在希腊也几乎如此，由罗马皇帝将\Person{安提诺乌斯}神化，他爱\Person{安提诺乌斯}正如\Person{宙斯}爱\Person{伽倪墨得斯}，其美貌极为罕见：因为情欲若无畏惧便难以约束；如今人们遵守\Person{安提诺乌斯}的神圣之夜，与他共度此夜的情人深知其可耻。为何将因不洁而受尊崇者列入诸神？为何命你哀悼他为子？为何你要夸耀他的美貌？被恶习玷污的美是可憎的。人啊，不要对美施行暴政，也不要对盛开的青春施加污辱。保持美的纯洁，使其真正美丽。作美之王，而非暴君。保持自由，那时我将承认你的美，因为你保持了其形象的纯洁：那时我将敬拜那真正的美，即一切美者的原型。如今这放荡少年的坟墓就是\Person{安提诺乌斯}的神庙和城镇。因为正如神庙受到敬畏，同样坟墓、金字塔、陵墓和迷宫也如此，它们是死者的神庙，如同前者是诸神的坟墓。关于这一点，我要向你们引证女先知\Person{西彼拉}的话：—\par

\Quote{非是\Person{福玻斯}的欺人神谕，\newline{}愚人称之为神，妄称为先知；\newline{}而是伟大天主的谕言，非人手所造，\newline{}如石刻的哑巴偶像。}\par

她也预言神庙的毁灭，预言说\Place{以弗所}的\Person{阿耳忒弥斯}神庙将被地震和地裂吞没，如下：—\par

\Quote{\Place{以弗所}将匍匐在地，在岸边哭泣哀号，\newline{}寻找一座再无居者的神庙。}\par

她还说\Person{伊西斯}和\Person{塞拉比斯}的神庙将被拆毁焚烧：—\par

\Quote{\Person{伊西斯}，三倍不幸的女神，你将逗留在\Place{尼罗河}之滨；\newline{}孤独、疯狂、沉默，在\Place{阿克戎}的沙滩上。}\par

接着她继续说道：—\par

\Quote{而你，\Person{塞拉比斯}，被一堆白石覆盖，\newline{}将躺在三倍不幸的\Place{埃及}，成为一片巨大的废墟。}\par

但若你们不听女先知，至少听听你们自己的哲学家，\Place{以弗所}的\Person{赫拉克利特}，他斥责偶像的无知觉：\Quote{他们向这些像祈祷，结果如同对自家墙壁说话。}因为那些敬拜石头、将其置于门前如同能行动者，岂不可惊？他们敬拜\Person{赫耳墨斯}为神，又立\Person{阿吉乌斯}为守门者。若有人责备它们无感觉，为何敬拜它们为神？若认为它们有感觉，又为何置于门前？罗马人将其最大成功归功于\Person{福尔图娜}，视其为极伟大的神祇，便将她的雕像搬进厕所，竖立在那里，为女神指定了一座合宜的神庙——那必需之所。但无感觉的木石和贵重的黄金，既不在乎香气，也不在乎血和烟，这些同时尊崇和熏染使它们变黑；同样也不在乎荣誉或侮辱。这些偶像比任何动物都更无价值。我无法理解无感觉之物如何被神化，并不得不怜悯那些迷失于这种愚昧迷宫中的可怜人：因为即使有些生物并非具备所有感官，如蠕虫和毛虫，以及那些甚至起初就不完全的，如鼹鼠和鼩鼱，\Person{尼坎德}说它盲目且粗鄙，它们仍优于那些完全无感觉的偶像和像。因为它们至少有一种感官——例如听觉、触觉，或类似嗅觉和味觉；而偶像连一种感官都没有。有许多生物既无视觉，也无听觉，也无言语，如牡蛎类，但它们仍然存活生长，并受月相变化影响。但偶像不动、不活、无感觉，被捆绑、钉住、粘合——被熔化、锉平、锯开、磨光、雕刻。无感觉的泥土被造像者玷污，他们用技艺改变其本性，诱导人敬拜它；造神者敬拜的不是神和鬼，依我之见，而是构成偶像的泥土和技艺。因为实际上，偶像只是匠人之手塑造的死物。但\Emph{我们}没有可感觉的物质形像，而是唯有心智可理解的形像——天主，唯有祂是真天主。\par

再者，当遭遇灾祸时，迷信的石头敬拜者虽已从事件中得知无感觉之物不应被敬拜，却因不幸的压力而成为其迷信的牺牲品；他们虽鄙视偶像，却又不愿显得完全忽视它们，便被那些以偶像之名称呼的神明所指责。\par

因为西西里的僭主小狄俄尼修斯从朱庇特雕像上剥下金披风，命令给他穿上羊毛的，并诙谐地说，后者比金的更好，夏天更轻便，冬天更暖和。库齐库斯的安条柯斯因缺钱，下令熔化十五肘高的金宙斯像，并制作一个类似但材质次等、包金的神像来代替。燕子与多数鸟类飞临这些雕像，在它们身上排泄，不顾奥林匹亚的宙斯、厄庇达乌里的阿斯克勒庇俄斯、甚至雅典娜·波利亚斯或埃及的塞拉比斯；但你们连从它们那里也没有学到偶像的愚蠢。而且，曾有歹徒或敌人袭击、焚烧神庙，劫掠奉献物，甚至因卑劣的贪欲熔化神像本身。如果有一个冈比西斯或大流士，或其他任何疯子这样做过，如果有人杀死了埃及的阿庇斯，我嘲笑他们杀死他们的神，同时为这出于贪欲的暴行感到痛苦。因此，我愿意忘记这样的恶行，视这些行为更多是贪婪之举，而非证明偶像的无能。但火与地震很精明，不会对魔鬼或偶像感到羞怯或害怕，正如对海浪堆在岸上的卵石一样。\par

我知道火能暴露并治愈迷信。若你愿意放弃这愚昧，火元素将为你的道路照明。同一把火烧毁了阿尔戈斯的神庙与女祭司克律西斯，以及以弗所的阿耳忒弥斯神庙（第二次，在阿玛宗人之后）。罗马的卡皮托利山多次起火；火也未放过亚历山大城中的塞拉比斯神庙。在雅典，它摧毁了厄琉忒里的狄俄倪索斯神庙；至于德尔斐的阿波罗神庙，首先遭受风暴侵袭，而后审慎之火将其彻底焚毁。这是火所应许的前言。造像者，难道他们不使你们中有智慧的人因鄙视物质而羞愧吗？雅典人菲狄亚斯在奥林匹亚的宙斯手指上刻下：\Quote{潘塔克斯是美的}。在他眼中，美的不是宙斯，而是他所爱的人。据波西狄普斯在其\WorkTitle{关于克尼多斯的书}中所载，普拉克西特列斯在制作克尼多斯的阿佛洛狄忒雕像时，将它造得像他所爱的克拉蒂纳的模样，使可怜的人们去敬拜普拉克西特列斯的情妇。当特斯皮亚的名妓弗吕涅风华正茂时，所有画家都将他们的阿佛洛狄忒画像摹写弗吕涅的美貌；同样，雅典的雕刻家使他们的墨丘利像酷似阿尔西比亚德。剩下由你们判断，是否该敬拜妓女。我相信，基于这些事实，并轻视这类荒诞故事，古代诸王毫无羞耻地宣称自己是神，因为这不会招致人的危险，从而教导说因他们的荣耀而成为不朽。埃俄罗斯之子刻宇克斯被其妻阿尔库俄涅称为宙斯；阿尔库俄涅又被其夫称为赫拉。托勒密四世被称为狄俄倪索斯；本都的米特拉达梯也被称为狄俄倪索斯；亚历山大希望被视为阿蒙之子，并让雕刻家为他造有角的雕像——渴望通过加角来玷污人形的美。不仅国王，连私人都以神名为荣，如医师墨涅克拉忒斯取了宙斯之名。我何需举阿勒克萨尔科斯为例？他原是文法学家，却扮作太阳神，据萨拉米斯的阿里斯托斯记载。何必提尼卡戈拉斯？他是【本都的】泽拉人，生活在亚历山大时代。尼卡戈拉斯自称赫尔墨斯，并穿赫尔墨斯的服饰，他自己为此作证。而整个民族和城市，全体居民沉溺于自谄，蔑视关于神的神话，同时人们自己，摆出与神平等的姿态，被虚妄的荣耀冲昏头脑，给自己施加过分的尊荣。有马其顿佩拉的腓力（阿明托尔之子）为例，他们在库诺萨戈斯为他颁定神明的崇拜，尽管他锁骨骨折、腿瘸、一只眼被击瞎。还有德米特里，被提升到神的行列；他在进入雅典时下马之处，便是德米特里 \Emph{下马者}的神庙，各地为他筑坛，雅典人将雅典娜嫁给他。但他蔑视女神，因他不能与雕像结婚；他带着名妓拉弥亚登上卫城，在雅典娜的床上与她同寝，向老处女展示年轻妓女的姿态。\par

因此，不必对希波感到愤慨，他使自己死亡不朽。这位希波下令在其墓碑上刻下如下的挽歌：\par

\Quote{这是希波之墓，命运\newline{}使他通过死亡与不朽之神等同。}\par

做得好，希波！你向我们展示了人的幻觉。如果他们不信活着的你说话，现在你死了，让他们成为你的门徒。这是希波的神谕，让我们思考。你们所崇拜的对象曾经是人，随时间而死亡；神话与时间使他们获得尊荣。因为某种方式，现实之物因熟悉而被轻视；而过去之物，因时间的模糊性而与临时的责难相隔，被虚构赋予尊荣，以至于现在被怀疑，过去被赞赏。正是这样，古时的死人因长久流行的谬误而受崇敬，被后代视为神。作为你们相信这些的依据，有你们的秘仪、庄严的集会、锁链与伤口，以及哭泣的神明。\par

\Quote{祸哉，祸哉！命运注定我最爱的萨尔珀冬\newline{}要被帕特洛克罗斯之手杀死。}\par

\Person{宙斯}的旨意被压倒了；\Person{宙斯}败北，他为\Person{萨耳珀冬}哀悼。因此，你们自己也合理地称他们为阴影和精灵，因为\Person{荷马}以不祥的尊荣敬奉\Person{雅典娜}和其他神祇，称他们为精灵：—\par

\Quote{她朝着天路飞行，\newline{}加入\Person{宙斯}居所中的不朽者。}\par

既然如此，阴影和精灵怎能仍被算作神祇？他们实际上是不洁不净的灵，众所公认属于土与水之性，因自身重量下沉，在坟墓和陵墓周围飘荡，隐约显现，不过是虚幻的幽灵罢了。你们的神祇就是这样的东西——阴影和影子；再加上那些残缺、皱缩、斜视的神祇\Person{利泰}，是\Person{忒耳西忒斯}的女儿，而非\Person{宙斯}的。因此，在我看来，\Person{比翁}机智地说：人怎能合理地向\Person{宙斯}祈求美丽的后裔——连他自己都未能得到的事呢？\par

那不朽的存在，就你们而言，你们把它沉入地下；那纯洁神圣的本质，你们将它埋葬在坟墓中，剥夺了神性的真实本性。\par

请问，你们为什么将天主的特权赋予那些并非神祇的东西？请问，你们为什么离弃天上，向大地致以神圣的尊荣？金子、银子、钢、铁、铜、象牙、宝石，又是什么呢？难道不都是泥土，出于泥土吗？\par

你们所看的这一切，不都是同一位母亲——大地——的后裔吗？\par

那么，愚蠢愚昧的人们（我要重复），你们为什么亵渎超天界领域，将宗教拖到地上，为自己塑造了泥土的神祇，并追随那些受造之物而非未受造的天主，沉入最深的黑暗之中？\par

\Place{帕罗斯}石很美，但它还不是\Person{波塞冬}。象牙很美，但它还不是\Place{奥林匹亚}的\Person{宙斯}。质料总需要艺术来塑造，但神性一无所缺。艺术前来做工，质料披上了形状；材料的珍贵使其能被利用来获利，但只是由于其形式，它才被认为值得敬拜。你们的偶像，若就其起源而言，是金子、木头、石头、泥土，从艺术家手中获得了形状。但我习惯行走在大地上，而非敬拜它。因为我将我的灵魂之望托付给缺乏生命气息的东西，是错误的。因此我们必须尽可能贴近偶像。欺骗如何内在地属于它们，从它们的外表就显而易见。因为偶像的形式明显带有精灵的本性印记。如果某人走一圈，检查图画和偶像，他瞥一眼就能从它们可耻的形式中认出你们的神祇：\Person{狄俄尼索斯}从他的袍子认出；\Person{赫准斯托斯}从他的技艺认出；\Person{得墨忒耳}从她的灾祸认出；\Person{伊诺}从她的头饰认出；\Person{波塞冬}从他的三叉戟认出；\Person{宙斯}从天鹅认出；火葬堆表明\Person{赫拉克勒斯}；如果某人看到一座裸体女人雕像而没有铭文，他明白那是金色的\Person{阿佛洛狄忒}。就这样，那个\Place{塞浦路斯}的\Person{皮格马利翁}爱上了一座象牙偶像：那偶像是\Person{阿佛洛狄忒}，并且是裸体的。塞浦路斯人仅凭形状就被征服，拥抱偶像。\Person{菲罗斯特法诺斯}讲述了此事。在\Place{奈多斯}有另一个不同的\Person{阿佛洛狄忒}，是石头的，而且美丽。另一个人爱上了它，可耻地拥抱那石头。\Person{波西狄波斯}讲述了此事。前者在其论塞浦路斯的书中，后者在其论\Place{奈多斯}的书中。艺术如此强大，足以欺骗，引诱好色之人陷入陷阱。艺术强大，但不能欺骗理性，也不能欺骗那些合乎理性生活的人。画中的鸽子被画家的技艺描绘得如此逼真，以致鸽子飞向它们；马匹曾对着画得逼真的母马嘶鸣。据说一个女孩爱上了一座偶像，一个俊美的青年爱上了\Place{奈多斯}的雕像。但观众的双眼被艺术欺骗了；因为没有一个神志清醒的人会拥抱一位女神，或与一个无生命的爱人一同埋葬，或爱上精灵和石头。但艺术以另一种咒语欺骗你们，如果它不引导你们沉溺于情欲：它引导你们向偶像和图画致以宗教尊崇和敬拜。\par

图画是像的。好极了！让艺术得到它应得的赞誉，但不要让艺术冒充真理而欺骗人。马安静地站着；鸽子不扑翅，翅膀不动。但\Person{代达罗斯}的木制母牛引诱了野公牛；艺术欺骗了它，迫使它遇到一个充满淫欲的女人。这就是恶作剧的艺术在无感觉者心中创造的狂热。另一方面，猿猴受到喂养照料它们的人的赞赏，因为没有什么像偶像和蜡或粘土制成的女孩饰物能欺骗它们。那么，你们将显得不如猿猴，因为你们依附于石头、木头、金子、象牙的偶像和图画。你们这些制造此类恶作剧玩具的人——雕刻家和偶像制造者、画家、金属工人、诗人——引入了一群混杂的神祇：田野里有\Person{萨提尔}和\Person{潘}；森林里有\Person{宁芙}、\Person{俄瑞阿得斯}和\Person{哈玛德律阿得斯}；此外，水中有河流、泉源的\Person{那伊阿得斯}；海中有\Person{涅瑞伊得斯}。如今，\OriginalTerm{玛吉}{Magi}吹嘘精灵是他们不敬的仆人，将他们算作家仆中的一员，并用咒语迫使他们做奴隶。此外，神祇的婚配、他们所叙述的生养儿女、歌中歌颂的通奸、喜剧中表现的纵酒宴乐、杯盏间的哄笑——你们的作者引入这些，促使我呼喊，尽管我宁愿沉默。哦，无神！你们把天上变成了舞台；神圣变成了戏剧；你们以精灵的面具在喜剧中扮演神圣的事物，以你们的\OriginalTerm{精灵崇拜}{demon-worship}（迷信）歪曲了真宗教。\par

\Quote{但他弹起竖琴，开始优美地歌唱。}\par

为我们歌唱，\Person{荷马}，那优美的歌\par

\Quote{关于\Person{阿瑞斯}与\Person{维纳斯}和美丽花冠的恋情：\newline{}起初他们怎样在\Person{赫菲斯托斯}的宫殿中秘密同床；\newline{}他送了许多礼物，玷污了\Person{赫菲斯托斯}王的床和房间。}\par

住口，哦，\Person{荷马}，别再唱那歌！它不美丽；它教导奸淫，我们被禁止听奸淫的故事以免玷污耳朵，因为我们是在这活生生且动态的人性形象中随身带着\OriginalTerm{天主的肖像}{likeness of God}——这一肖像与我们同住，与我们商议，与我们交往，作客于我们，与我们同感，为我们感同身受。我们因\Person{基督}的缘故成了\OriginalTerm{祝圣的奉献}{consecrated offering}归于天主：我们是\OriginalTerm{特选的种族}{chosen generation}，\OriginalTerm{王家的司祭}{royal priesthood}，\OriginalTerm{圣洁的国民}{holy nation}，\OriginalTerm{属主的子民}{peculiar people}，从前不是子民，如今却是天主的子民；按照\Person{若望}的话，我们不属于那些在下的人，而是从那位从上而来的学习了万有；我们理解了天主的救恩计划；我们学会了在生命的更新中行走。但这并不是多数人的想法；相反，他们抛弃羞耻和恐惧，在家里描绘魔鬼的不洁情欲。因此，他们与污秽结盟，用挂在卧室中的彩绘板装饰自己的卧室，将放纵视为宗教；躺在床上的拥吻之际，他们观看那\Person{阿佛洛狄忒}与她的情人相拥。他们在戒指的环上雕刻一只绕着\Person{勒达}盘旋的情鸟——对柔弱的描绘有强烈偏好——并使用的印章上印着\Person{宙斯}放纵的印记。这些便是你们淫逸的例证，这些便是邪恶的神学，这些便是你们与你们一同行淫的神祇的教导；因为正如那位雅典演说家所说，人想什么，就信什么。另一方面，你们其他的偶像又是怎样的呢？小巧的\Person{潘}神，裸体的少女，醉醺醺的\Person{萨提尔}，以及象征阳具的标记，在画中赤身露体，污秽不堪。不仅如此：你们不羞于在众人眼前观看公共场所描绘的所有形式的放荡图像，反而将它们陈列出来并精心守护，将这些无耻的柱子在家中奉为神圣，仿佛它们真是你们的神祇的像，在上面同样描绘着\Person{斐莱尼斯}的姿态和\Person{赫拉克勒斯}的苦役。我们不仅谴责这些的使用，连看到和听到它们，也认为该受遗忘的厄运。你们的耳朵被败坏，你们的眼睛行淫，你们的眼神在拥抱之前就已犯奸淫。哦，你们这些对人施暴并将这\Person{天主}手造物中的神圣之物献于羞耻的人，你们不信一切，以便放纵情欲，并且相信偶像，因为你们渴望它们的放荡，却不相信天主，因为你们无法忍受自律的生活。你们恨更好的，珍视更坏的，你们是德性的旁观者，却是恶行的实行者。因此，可以说，所有那些人都一同有福——\par

\Quote{谁拒绝看任何庙宇\newline{}和祭台——那些静默石头的无价值坐席，\newline{}以及石雕的偶像和人手制作的像，\newline{}沾染了活物的血和牺牲\newline{}——四足兽、两足兽、飞禽和野兽的屠宰。}\par

因为我们明确被禁止行欺骗之术：先知说：\BibleQuote{你不可制造任何像天上或地上的事物之像。}\BibleRef{出 20:4}\par

难道我们还能再以为\Person{普拉克西特列斯}的\Person{得墨特尔}、\Person{科瑞}和神秘的\Person{伊阿科斯}是神，而不认为\Person{留基伯}的艺术或\Person{阿佩莱斯}的手——它们以神圣荣耀的形象装扮物质——更配得尊敬吗？当你们竭尽全力使偶像尽可能精美地成形时，却毫不留心避免变得像偶像一样愚钝。因此，先知的话以最清晰简洁的方式谴责这种行为：\BibleQuote{因为列邦的所有神祇都是魔鬼的像，但天主创造了诸天和其中的万有。}\BibleRef{咏 96:5}\newline{}有些人却不知何故陷入错误，敬拜天主的工程而非天主本身——太阳、月亮和其余星辰的唱诗班——荒谬地认为这些只是测量时间的工具是神；\BibleQuote{因为祂的话使他们确立，祂口中的气使他们众多。}\BibleRef{咏 33:6}\par

人的艺术进一步产生房屋、船只、城市和图画。但我怎能述说天主所造的？看整个宇宙：这是祂的工程；诸天、太阳、天使和人都出自祂的手指。天主的权能何其伟大！祂单单的旨意便是宇宙的创造。因为只有天主创造了它，因为只有祂是真正的天主。祂仅仅通过运用旨意而创造；祂的一念便使祂所意愿的立刻成真。因此，哲学家的合唱队错了，他们确实最高贵地承认人被造是为了默观诸天，却敬拜诸天中显现且被眼睛捕捉的对象。因为若天体不是人的作品，它们当然是为人类造的。你们中谁也不要敬拜太阳，却当渴慕太阳的创造者；也不要神化宇宙，而当寻求宇宙的创造者。那么，想进入救恩之门的人唯一剩下的避难所便是\OriginalTerm{天主的智慧}{divine wisdom}。由此，如同从神圣的避难所，追求救恩的人不能被任何魔鬼拖走。\par

\SourceRecord{Translated by William Wilson.；Ante-Nicene Fathers；Vol. 2.；Edited by Alexander Roberts, James Donaldson, and A. Cleveland Coxe.；Buffalo, NY: Christian Literature Publishing Co.,；1885.；<http://www.newadvent.org/fathers/020804.htm>.}

% structure: p0001 source=h1 target=section reason=page-title

\section{劝告外邦人（第五章）}

\Person{亚历山大的克莱孟} 著\par

% structure: p0003 source=h2 target=subsection reason=relative-heading-rank; base=h2

\subsection{第五章 哲学家们关于天主的意见}

那么，若你愿意，让我们来检视一下哲学家们对诸神所夸耀的意见；这样我们就可以发现，哲学本身因傲慢而将物质当作偶像崇拜，尽管我们能够随着论述的展开而指出，它甚至在神化某些魔鬼时，也怀有对真理的梦想。他们中有些人将元素指定为万物的本原：米利都的泰勒斯歌颂水，同样来自米利都的阿那克西美尼则以空气为万物的本原，后来阿波罗尼亚的第欧根尼追随了他。埃利亚的巴门尼德引入了火和土作为神；其中火这一元素，梅塔蓬图的希帕索斯和以弗所的赫拉克利特都认为它是神明。阿格里根特的恩培多克勒则追随了众多神明，除了那四种元素之外，还列举了不和与友爱。这些人无疑应被视为无神论者，他们通过不明智的智慧崇拜物质，虽未向树木和石头表示宗教上的敬意，却神化了大地——这些物质的母亲；他们没有为波塞冬制作雕像，却崇敬水本身。因为波塞冬根据其原初含义，若非是潮湿的物质，又是什么呢？其名称源自 \OriginalTerm{饮用}{posis}；同样，好战的阿瑞斯之所以如此称呼，无疑源自 \OriginalTerm{上升}{arsis} 和 \OriginalTerm{毁灭}{anœresis}。我认为主要是这个原因，许多人将一把剑插在地上，当作阿瑞斯来献祭。西徐亚人就有这种习俗，正如欧多克索斯在其 \WorkTitle{游记} 第二卷中所告诉我们的。而萨罗马泰人，一支西徐亚人的部落，也崇拜一把军刀，一如伊克西乌斯在其 \WorkTitle{奥秘} 中所说。\par

赫拉克利特及其追随者也是如此，他们崇拜火为第一因；这火被其他人称为赫淮斯托斯。波斯祆教徒以及许多亚洲居民也崇拜火；除此之外，马其顿人也是如此，正如第欧根尼在其 \WorkTitle{波斯志} 第一卷中所记载的。为何要特别提及萨罗马泰人呢？尼姆福多鲁斯在其 \WorkTitle{蛮族风俗} 中说他们尊崇火为神圣之物。至于波斯人、米底人或祆教徒呢？狄诺告诉我们，他们在露天之下献祭，视火与水为神祇的唯一形象。\par

我也没有隐瞒他们的无知；因为，无论他们多么自以为以一种形式避免了错误，却以另一种形式滑入了错误。\par

他们不像希腊人那样，认为树木和石头是神祇的像；也不像埃及人那样，认为朱鹭和猫鼬是神祇的像；而是像哲学家那样，认为火和水是神祇的像。贝罗苏斯在其 \WorkTitle{迦勒底史} 第三卷中指出，人类是在许多个连续的年代之后才开始崇拜人形偶像的，这一习俗是由阿塔薛西斯——大流士之子、奥库斯之父——引入的，他首先在巴比伦和苏萨设立了阿佛洛狄忒·阿纳伊提斯的偶像；而埃克巴塔那则向波斯人作出了崇拜它的榜样；巴克特里亚人则向大马士革和萨迪斯作出了榜样。\par

那么，让哲学家们承认波斯人或萨罗马泰人或祆教徒是他们的老师吧，他们从这些人那里学到了将某些本原视为神明的邪恶教义，却不认识那伟大的第一因，万物的创造者，就连那些本原本身的创造者，无始的天主，反而敬畏保禄所说的\BibleQuote{这些懦弱无用的基本元素}\BibleRef{迦 4:9}，它们本是为服务于人而被造的。至于其余的那些超越了元素、热切探求更高更崇高之物的哲学家们，有些人大谈无限，如米利都的阿那克西曼德、克拉佐美纳的阿那克萨戈拉和雅典的阿凯劳斯，这两人都将 \OriginalTerm{理智}{\GreekText{νοῦς}} 置于无限之上；而米利都的留基伯和希俄斯的梅特罗多洛似乎教导了两个本原——充满与虚空。阿布德拉的德谟克利特虽然接受了这两者，又添加了 \OriginalTerm{影像}{\GreekText{εἴδωλα}}；克罗托内的阿尔克迈翁则主张星辰是神祇，且拥有生命（对于他们的厚颜无耻，我缄口不言）。卡尔西顿的色诺克拉底指出七颗行星是七位神明，而由这一切构成的宇宙是第八位。我也不会放过廊下派的人，他们说神性渗透于一切物质，甚至最卑贱的物质之中，从而笨拙地玷污了哲学。我认为，既然已经谈到这一步，顺便提及逍遥派也不算冒犯。这一学派的创始人，不认识万有的父，认为那被称为至高者的是宇宙的灵魂；也就是说，他假设宇宙的灵魂就是神，因而被自己的剑刺穿了。因为他首先把天命的范围限定到月球轨道，然后又假定宇宙是神，从而自相矛盾，因为他教导说，那没有神的宇宙却是神。而那位埃雷索斯的泰奥弗拉斯托斯，亚里士多德的门徒，有时猜测天是神，有时猜测精神是神。唯有伊壁鸠鲁，我很乐意忘记他，他将不敬推到极致，认为天主毫不关心世界。此外，本都的赫拉克利德呢？他被到处拖向德谟克利特的\OriginalTerm{影像}{\GreekText{εἴδωλα}}。\par

\SourceRecord{Translated by William Wilson.；Ante-Nicene Fathers；Vol. 2.；Edited by Alexander Roberts, James Donaldson, and A. Cleveland Coxe.；Buffalo, NY: Christian Literature Publishing Co.,；1885.；<http://www.newadvent.org/fathers/020805.htm>.}

% structure: p0001 source=h1 target=section reason=page-title

\section{对异教徒的劝勉（第六章）}

\Signature{\Person{亚力山大的克肋孟}}\par

% structure: p0003 source=h2 target=subsection reason=relative-heading-rank; base=h2

\subsection{第六章 因着天主灵感，哲人有时触及真理。}

一大批这样的思绪涌上我的心头，仿佛带来了狰狞可怖的异形鬼怪，以老妇之言谈论着神话般的奇形怪状。我们绝非劝人听信此类故事，我们避免像俗语所说，用荒唐故事哄哭闹的孩子，因为我们担心在他们心中助长那些自以为是、其实对真理一无所知如婴孩之人的不敬。因为（以真理之名！）你们为何使那些相信你们的人堕入毁灭与腐朽，可怕而不可挽回？我恳求你们，为何用偶像的形像充满生命，妄称风、气、火、土、石头、木头、钢铁，或这宇宙为神，并在那备受赞誉的占星术（而非天文学）中，向那些真正迷失之人夸耀天体，称流浪之星为神？我所渴求的是万灵之主、火之主、宇宙的创造者、那点燃太阳者。我寻求的是天主，而非天主的工程。我当寻求谁为助伴？你若无异议，我们并非完全摒弃\Person{柏拉图}。那么，\Person{柏拉图}，如何探寻天主？\Quote{因为要找到这宇宙的父和造物主是件困难的事；找到后，完全宣示祂是不可能的。}\par

为何如此？我恳求你，凭祂自身！因为祂绝不能被表达。说得好，\Person{柏拉图}！你触及了真理。但不要懈怠。与我一同探究至善。因为某种神圣的流溢已注入所有人，尤其是那些专注于理智追求的人；因此，他们虽不情愿，也承认天主是唯一的、不可毁灭的、非受生的，并且在天上的某处，在祂自己独特而恰当的高处，祂俯瞰万物，拥有真实而永恒的存在。\par

\Quote{请告诉我，我该将天主想作什么，\newline{}祂看见一切，而自身却不可见。}\newline{}\Person{欧里庇得斯}如此说。\par

据此，\Person{梅南德}在我看来似乎犯了错误，当他说：——\par

\Quote{哦，太阳！因为你，诸神中的首位，应当受敬拜，\newline{}因着你的恩赐，我们得以看见其他神祇。}\par

因为太阳永远不能给我显示真天主；唯有那健康的圣言——灵魂的太阳——当它在灵魂深处升起时，灵魂之眼本身才被光照。因此，\Person{德谟克利特}不无道理地说：\Quote{少数有才智的人，将手举向上方——我们希腊人如今称之为\OriginalTerm{空气}{\GreekText{ἀήρ}}——把整个空间称作宙斯或天主；祂也知晓一切，给予并收回，祂是万王之王。}\par

\Person{柏拉图}也有同样的见解，他在某处如此暗示天主：\Quote{万物都围绕着万王之王，祂是一切美善事物的原因。}那么，万王之王是谁？天主，祂是一切存在之真理的尺度。因此，正如被度量的事物包含在尺度中，对天主的认知也度量和包含真理。真正神圣的\Person{梅瑟}说：\BibleQuote{在你袋中不可有一大一小两样的法码，但须有真实公平的法码。}他认为整全的尺度、度量、数目就是天主。因为不义不洁的偶像藏在袋中，即藏在污秽的灵魂里。唯独公平的尺度是唯一真实的天主，永远公义，恒久不变；祂度量一切，以正义如天平般称量，把握并维持整个本性于平衡之中。\Quote{因此，天主，如古语所说，占据一切存在者的始、中、终，沿着直路运行，同时环行自然；正义永远跟随祂，惩罚那些违反神圣律法的人。}\par

\Person{柏拉图}啊，你从何处得到那真理的暗示？何处得到这丰富雄辩的言辞，以神谕的方式宣告虔诚？他说，野蛮人的部族比这些更智慧；我知道你的老师，即使你隐瞒。你从埃及人学到几何，从巴比伦人学到天文；从色雷斯人获得治疗的秘术；亚述人也教了你许多；但关于符合真理的法律和关于天主的见解，你归功于希伯来人。\par

\Quote{他们不因虚妄的欺骗敬拜\newline{}人手所造之物，金、铜、银、象牙，\newline{}和死人的像，木、石之偶，\newline{}这些被其他人因愚昧倾向所敬拜；\newline{}却向天举起洁净的双臂：\newline{}他们起床时用水洁净自己，\newline{}唯独敬拜永恒者，祂统治直到永远。}\par

且不止\Person{柏拉图}一人；哦，哲学啊，快带出许多其他人来，他们凭祂的启示宣告唯一的真天主是天主，只要他们多少抓住了真理。因为\Person{安提斯泰尼}并非自己想出犬儒学说；而是因他身为\Person{苏格拉底}的门徒，他说：\Quote{天主不像任何事物；因此无人能从形像认识祂。}雅典的\Person{色诺芬}本人本可自由写下某些真理，并作与\Person{苏格拉底}同样的见证，若非他惧怕\Person{苏格拉底}所饮的毒酒。但他仍暗示了同样的事：他说：\Quote{那推动万物而自身静止的天主何其伟大有力，显然可见；但祂的形态如何却未显明。太阳本身，原要为周围一切提供光源，却不认为适宜让人注视；若有人胆敢凝视它，便失去视力。}那么，格里鲁斯的儿子从何处学到他的智慧？岂不显然来自那位希伯来女先知，她以下列方式预言：——\par

\Quote{什么血肉之躯能用眼睛看见\newline{}那属天的、真实不朽的天主，居住在天穹之上？\newline{}不，生于死亡的世人甚至不能\newline{}站立在太阳的光芒之前。}\par

\Person{克莱安忒斯·皮萨迪乌斯}，斯多葛派哲学家，他所展露的并非诗意的神谱，而是真实的神学，毫不隐藏自己对天主所持的看法：——\par

\Quote{你若问我善的本性为何，且听：\newline{}就是那规整、公正、神圣、虔诚之物，\newline{}自律、有益、公平、适宜，\newline{}庄重、独立、常有益处；\newline{}那物无惧无忧；有益、无痛，\newline{}助人、愉快、安全、友善；\newline{}受尊敬、与自身一致、可敬；\newline{}谦卑、谨慎、温良、热心，\newline{}持久、无瑕、永远长存：\newline{}凡着眼于意见的人便是卑微，\newline{}为从中获取一些利益。}\par

这里，依我之见，他清楚地教导了天主的本性；并且指出，世俗的见解与宗教习俗奴役那些追随它们却不寻求天主的人。\par

我们也不可将毕达哥拉斯学派置于幕后，他们说：\Quote{天主是唯一的；他并非像某些人所想的那样，存在于宇宙框架之外，而是内在于其中；且以其整个存有的一切完满，存在于整个存在界的循环中，审察整个本性，并将万有和谐地融合为一——他是自己所有力量与工程的创造者，天上光明的赐予者，万有之父——整个世界的理智与活力——万物的推动者。}因为关于天主的认识，这些由我们提到的那些人受天主默感而写下的、并经我们选录的言语，甚至对那只有少许能力探究真理的人也足够了。\par

\SourceRecord{Translated by William Wilson.；Ante-Nicene Fathers；Vol. 2.；Edited by Alexander Roberts, James Donaldson, and A. Cleveland Coxe.；Buffalo, NY: Christian Literature Publishing Co.,；1885.；<http://www.newadvent.org/fathers/020806.htm>.}

% structure: p0001 source=h1 target=section reason=page-title

\section{劝勉外邦人（第七章）}

\Signature{克莱孟（亚历山大城）}\par

% structure: p0003 source=h2 target=subsection reason=relative-heading-rank; base=h2

\subsection{第七章　诗人们也为真理作证}

让诗歌也来接近我们（因为仅凭哲学是不够的）；诗歌虽全然充斥着虚假，却几乎不愿承认真理，反倒会向天主承认自己偏离到神话中。让那些诗人中任何一位首先上前吧。\newline{}\Person{阿拉托斯}认为天主的权能渗透万有：——\par

\Quote{为使一切安稳，\newline{}他们始终首先并最后祈求祂，\newline{}万福，圣父啊！伟大的奇迹，对人的巨大恩赐。}\par

同样，阿斯克拉的\Person{赫西俄德}也隐约谈及天主：——\par

\Quote{因为祂是万有的君王，\newline{}不死者的主宰；没有谁能\newline{}与祂的权能抗衡。}\par

在舞台上，他们也揭示真理：——\par

\Quote{仰望以太和天，视其为天主，}\par

——\Person{欧里庇得斯}说。而\Person{索福克勒斯}，\Person{索菲洛斯}之子，说：——\par

\Quote{一，确实，仅有一位天主，\newline{}祂造了天和广袤的地，\newline{}以及海洋的碧波和强劲的风；\newline{}但我们许多有死之人，心受蒙蔽，\newline{}为自己立起了——作为逆境中的慰藉——\newline{}石头、木头、青铜、\newline{}黄金或象牙所造的神像；\newline{}并为那些神像献祭和举行徒然的节庆集会，\newline{}习惯于如此实践宗教。}\par

他就这样大胆地在舞台上将真理摆在观众面前。但色雷斯人\Person{俄耳甫斯}，\Person{俄阿格洛斯}之子，既是秘仪祭司又是诗人，在描述秘仪和偶像神学之后，以真正的庄严引入了真理的翻案诗，虽然姗姗来迟地唱道：——\par

\Quote{我将向那些蒙允许的人宣说；但要关上门，\newline{}无论如何，反对一切不敬者。但你听，\newline{}哦\Person{穆赛俄斯}，光明的月亮所生之子，\newline{}因为我要宣告何为真实。不要让这些\newline{}曾出现在你胸中的事物夺去你宝贵的生命；\newline{}但要仰望神圣的言语，专注于它，\newline{}保持\Emph{理智和情感的座位}正直，并行走在\newline{}笔直的道路上，只向宇宙的不朽君王\newline{}定睛注视。}\par

接着，他继续清楚地说：——\par

\Quote{祂是独一的，自存的；万物从祂独一而来，\newline{}祂自身在万物中活动：没有凡人\newline{}看见祂，但祂看见万有。}\par

至此\Person{俄耳甫斯}终于明白自己犯了错误：——\par

\Quote{但不要再犹豫了，哦，富有多种智慧的人；\newline{}回转吧，退返你的脚步，求悦于天主。}\par

因为，最多而言，希腊人领受了神圣言语的些许火花，发出了某些真实的话语；他们确实见证真理的力量并不隐藏，同时也暴露了自己的软弱——未能抵达终点。因为我认为现在对所有人已经清楚：那些没有真理之言而行事或说话的人，就像被迫无脚行走之人。愿诗人们因真理之力而在喜剧中引入的对你们诸神的斥责，使你们羞愧而得到救恩。例如，喜剧诗人\Person{米南德}在他的剧作\WorkTitle{驭手}中说：——\par

\Quote{没有一个神令我喜悦，他\newline{}跟着一个老妇人，进入房屋，\newline{}端着一个盘子。}\par

因为\Person{库柏勒}的乞讨祭司就是如此。因此\Person{安提斯泰尼}对他们的乞讨作了恰当的回答：——\par

\Quote{我不供养众神之母，\newline{}因为众神供养她。}\par

同样，那位喜剧作家不满意通常的习俗，试图在剧作\WorkTitle{女祭司}中揭露流行错误的邪恶傲慢，明智地宣告：——\par

\Quote{如果一个人用手鼓的响声\newline{}将神祇随意拖曳，\newline{}做此事的人就比神祇更大；\newline{}但这些乃是狂妄的工具和谋生的手段，\newline{}由人发明。}\par

不仅\Person{米南德}，而且\Person{荷马}、\Person{欧里庇得斯}以及其他大量诗人，都揭露你们的诸神，并惯于斥责他们，而且斥责得相当厉害。例如，他们称\Person{阿佛洛狄忒}为狗蝇，称\Person{赫淮斯托斯}为瘸子。\Person{海伦}对\Person{阿佛洛狄忒}说：——\par

\Quote{你的神位退位吧！\newline{}放弃奥林波斯！}\par

关于\Person{狄俄尼索斯}，\Person{荷马}毫无保留地写道：——\par

\Quote{他在可爱的\Place{倪萨}的树林中，\newline{}{在}他们的狂欢秘仪里，\newline{}使年轻的\Person{巴科斯}的乳母们狼狈溃逃；\newline{}她们恐惧，\newline{}各自丢弃了神杖，被凶猛的\Person{吕枯尔戈斯}手持刺棒驱散。}\par

真正配得上苏格拉底学派的是\Person{欧里庇得斯}，他注目真理，藐视看戏的观众。有一次，\Person{阿波罗}，\par

\Quote{祂居住在大地中央的圣所，\newline{}向凡人所发最确切的神谕，}\par

就这样被揭露：——\par

\Quote{我服从他杀了生我的她；\newline{}你视他为有罪——将他处死；\newline{}是他犯了罪，而不是我，因为我未受教导\newline{}{关于}正义与公平。}\par

他引入\Person{赫拉克勒斯}时而疯狂，时而醉酒贪食。他岂能不如此描绘那位神祇？当神祇被招待为客时，他大嚼生绿无花果，发出刺耳的嚎叫，甚至连他那蛮族的东道主都注意到了。在他的剧作\WorkTitle{伊昂}中，他也厚颜无耻地将诸神搬上舞台：——\par

\Quote{那么，你们这些为凡人立法者，\newline{}自己却行不义，这如何合理？\newline{}如果——这绝不会发生，但我仍如此假设——\newline{}你们要向人交出你们通奸的帐，\newline{}你，\Person{波塞冬}，还有你，\Person{宙斯}，天界的统治者，——\newline{}你们将为了补偿你们的恶行\newline{}不得不腾空你们的圣殿。}\par

\SourceRecord{Translated by William Wilson.；Ante-Nicene Fathers；Vol. 2.；Edited by Alexander Roberts, James Donaldson, and A. Cleveland Coxe.；Buffalo, NY: Christian Literature Publishing Co.,；1885.；<http://www.newadvent.org/fathers/020807.htm>.}

% structure: p0001 source=h1 target=section reason=page-title

\section{劝勉外邦人（第八章）}

\Signature{克莱孟·亚历山大}\par

% structure: p0003 source=h2 target=subsection reason=relative-heading-rank; base=h2

\subsection{第八章 应在先知书中寻求真道}

现在，我们已经依次处理了其他要点，是时候进入先知书了；因为神谕为我们提供了获得虔敬所需的工具，从而确立了真理。神圣的圣经与智慧典籍构成了救恩的捷径。它们没有修饰，没有外在的言语之美，没有浮夸和诱惑，却高举被罪恶窒息的人类，教导人们藐视人生的变故；并用同一个声音医治许多邪恶，同时劝诫我们远离有害的欺骗，并清楚地劝勉我们获得摆在我们面前的救恩。那么，就让女先知西卜尔首先为我们唱出救恩之歌：——\par

\Quote{因此祂是完全可靠、无误的：\newline{}来，不再跟随黑暗和幽暗；\newline{}看，太阳那甜美的光芒荣耀照耀。\newline{}认识智慧，并将其珍藏在心中：\newline{}有一位天主，祂降下雨水、风、地震、\newline{}雷电、饥荒、瘟疫和悲惨的忧伤，\newline{}以及雪和冰。但为何详述细节？\newline{}祂统治上天，祂治理大地，\newline{}祂真实存在；}——\par

她与灵感惊人地一致，将错觉比作黑暗，将认识天主比作太阳和光明，并将两者加以比较，表明我们应当作出的选择。因为虚假并非因真理的单纯显示而消散，而是通过真理的实践运用，将其逐出并驱散。\par

先知\Person{耶肋米亚}，赋有卓越的智慧，更确切地说，是耶肋米亚内的圣神，彰显了天主。\BibleQuote{我是近处的天主，} 他说，\BibleQuote{不是远处的天主吗？人在隐密处行事，我岂看不见吗？我岂不充满天地吗？主说。} \BibleRef{耶 23:23}\par

又通过依撒意亚说：\BibleQuote{谁曾用手量过苍天，用掌测过大地？}\BibleRef{依 40:12}请看天主的伟大，并充满惊异。让我们朝拜祂，先知论祂说：\BibleQuote{在祢面前，山岭必熔化，如蜡在火前熔化！}\BibleRef{依 64:1}他说，这是\BibleQuote{天主的宝座是天，祂的脚凳是地；若祂敞开天，战栗必将你抓住。}\BibleRef{依 66:1}你还要听这位先知论偶像的话吗？\Quote{他们必在太阳面前成为示众之物，他们的尸首必成为天上飞鸟和地上野兽的食物；他们必在他们所爱所事奉的太阳和月亮面前腐烂；他们的城也必被焚毁。}他又说，元素和世界必被毁灭。\Quote{大地，}他说，\Quote{必变老，天必过去；但主的话永远长存。}那么，当天主再次藉梅瑟显示自己时：\BibleQuote{看，看，那是我，除我以外没有别的神。我杀人，也使人活；我击打，也医治；无人能从我手中救出。}\BibleRef{申 32:39}但你还愿听另一位先知吗？你有整个先知唱诗班，梅瑟的同工。圣神藉欧瑟亚所说的，我必不避讳引用：\BibleQuote{看，是我设立雷霆，创造神灵；祂的手建立了天上的万军。}\BibleRef{亚 4:13}又藉依撒意亚。我要重复这句话：\BibleQuote{我是，}他说，\BibleQuote{我是上主，我讲正义，宣告真理。你们聚集前来商议吧，你们这些从列国得救的人。那些立木块、雕刻物，并向不能救人的神祈祷的人，他们一无所知。}\BibleRef{依 45:19}接着：\BibleQuote{我是天主，除我以外没有公义的天主和救主；除我以外没有别的。你们这地极的人，都归向我，就必得救。我是天主，没有别的；我凭自己起誓。}\BibleRef{依 45:21}但祂对崇拜偶像的人发怒，说：\BibleQuote{你们将上主比作谁，或用什么像与祂相比？工匠岂非做了像，金匠熔金镀金？}\BibleRef{依 40:18}——等等。所以你们不要成为偶像崇拜者，甚至现在要警惕那些威胁；\BibleQuote{因为雕刻的偶像和人手的作品必哀号，更确切说，是那些倚靠它们的人必哀号，}\BibleRef{依 10:10}因为物质没有感觉。他又一次说：\BibleQuote{上主必震动居住的城邑，用手抓世界如抓巢穴。}\BibleRef{依 10:14}我何必向你们重复智慧的奥秘，以及希伯来人之子、智慧大师的著作中的话？\BibleQuote{上主造了我作为祂道路的起始，为祂的作为。}\BibleRef{箴 8:22}\Quote{上主赐予智慧，从祂面前发出知识和明达。}\BibleRef{箴 2:6}\BibleQuote{懒汉，你躺卧到何时？你何时从睡眠中醒来？}\BibleRef{箴 6:9}\BibleQuote{但若你显出不懒散，你的收获必如泉源涌来，}\BibleRef{箴 6:11}\BibleQuote{圣父的圣言，慈祥的光，带来光明的主，给众人的信德和救恩。}\BibleRef{箴 6:23}因为\BibleQuote{上主以自己的能力创造了大地，}正如耶肋米亚所说，\BibleQuote{以自己的智慧建立了世界，}\BibleRef{耶 10:12}因为智慧，即祂的圣言，将我们这些曾俯拜偶像的人举起归向真理，并且祂本身就是我们从堕落中的首次复活。因此，天主的人梅瑟劝诫杜绝一切偶像崇拜，美哉呼喊道：\Quote{以色列啊，请听，上主你的天主是唯一的主；你当朝拜上主你的天主，惟独事奉祂。}\Quote{所以，人哪，现今要明智，}按那有福的圣咏作者达味的话；\Quote{持守训诲，免得上主发怒，你们便从义路灭亡，因祂的怒气快要发作。凡投靠祂的人有福了。}但上主已以超然的怜悯，激发了救恩之歌，如同战歌：\Quote{人子啊，你们心迟慢要到几时？为何喜爱虚妄，寻求谎言？}那么，什么是虚妄？什么是谎言？主的圣宗徒责备希腊人，会告诉你们：\Quote{因为他们虽然认识天主，却没有当作天主光荣祂，也不感谢；反而在思想中成为虚妄的，将天主的荣耀改为可朽坏人的样式，敬拜事奉受造物多于造物主。}确实，这就是那位\BibleQuote{起初创造了天和地}\BibleRef{创 1:1}的天主。但你们不认识天主，却敬拜天，怎能逃脱不敬之罪？再听先知的话：\Quote{太阳必受亏损，天也昏暗；但全能者永远发光。那时，天上的万军将震动，天被伸展、卷起，如皮卷（这是先知的话），大地必从主面前逃走。}\par

\SourceRecord{Translated by William Wilson.；Ante-Nicene Fathers；Vol. 2.；Edited by Alexander Roberts, James Donaldson, and A. Cleveland Coxe.；Buffalo, NY: Christian Literature Publishing Co.,；1885.；<http://www.newadvent.org/fathers/020808.htm>.}

% structure: p0001 source=h1 target=section reason=page-title

\section{劝勉外邦人（第九章）}

\Signature{克莱孟（亚历山大的）}\par

% structure: p0003 source=h2 target=subsection reason=relative-heading-rank; base=h2

\subsection{第九章　\Quote{凡轻视或忽视天主慈恩召唤者，实犯严重犯罪}}

我能够列举成千上万的经文，没有一句\BibleQuote{一笔一划会落空}\BibleRef{玛 5:18}而不应验的，因为主圣神的口说了这些事。祂说：\BibleQuote{我儿，不要再轻视主的管教，被祂责备时也不要灰心。}\BibleRef{箴 3:11} 超越人情的爱！主不是像教师对学生、主人对仆人、或天主对人说话，而是像父亲温柔地劝告自己的孩子。梅瑟承认，当他听天主谈论圣言时，\BibleQuote{充满了战栗和恐惧}\BibleRef{希 12:21}。你听到天主圣言的声音，难道不害怕吗？你不感到痛苦吗？你不恐惧，急忙向祂学习——即得救——畏惧义怒、热爱恩宠、热切追求摆在我们面前的希望，好逃避所警告的审判吗？来吧，来吧，我的年轻人！因为如果你们不变成小孩子，不重生，正如圣经所说，你们不会接受那真实存在的父，也绝不能进入天国。因为外人如何被允许进入？我想，当他被登记并成为公民，并接受一个人作他的父亲时，他就会关心父亲的事，就会被认为配作祂的继承人，就会与祂的爱子分享父的国度。因为这是首生的教会，由许多好儿女组成；这些人就是\BibleQuote{登记在天上的首生者，与千千万万的天使一同欢庆。}我们也是首生的儿子，由天主养育，是首生者的真诚朋友，最先认识天主，最先脱离我们的罪恶，最先与魔鬼分离。如今，天主越是仁慈，人越是邪恶；因为祂希望我们从奴隶变成儿子，而他们却轻视成为儿子。以主为耻是何等愚昧！祂赐自由，你们却逃向奴役；祂赐救恩，你们却沉沦于毁灭；祂赐永生，你们却等待惩罚，宁愿选择主\BibleQuote{为魔鬼及其使者预备的}火。\BibleRef{玛 25:41} 因此，蒙福的宗徒说：\BibleQuote{我在主内作证：你们不要再像外邦人一样，存着虚妄的心思行事；他们的理智昏暗，因心中的顽固而与天主的生命隔绝，由于无知；他们既已麻木，便纵情淫欲，行各样不洁和贪欲。}\BibleRef{弗 4:17} 在如此见证的控告和宗徒对天主的呼求之后，对于不信者除了审判和定罪，还有什么呢？主不断地劝勉、恐吓、敦促、唤醒、告诫；祂从黑暗的睡梦中唤醒，并扶起那些在错误中徘徊的人。祂说：\BibleQuote{你这睡着的人，醒来吧，从死者中起来，基督必要光照你}\BibleRef{弗 5:14}——基督，复活的太阳，祂在晨星之前出生，以祂的光辉赋予生命。因此，不要有人轻视圣言，免得无意中轻视自己。因为经上某处说：\BibleQuote{今天，如果你们听从祂的声音，不要心硬，像在激怒时，在旷野试探的日子，你们的祖先试探了我。}那试探是什么？如果你想知道，圣神会指示你：\BibleQuote{并且看见了我的作为，四十年之久。所以我厌烦那一世代，说：他们心里总是迷误，不认识我的道路。于是我在忿怒中起誓：他们绝不得进入我的安息。}看看这警告！看看这劝勉！看看这惩罚！那么，我们为何还把恩宠变成忿怒，不肯侧耳聆听圣言，以纯洁的灵接待天主为客？因为祂的应许的恩宠何等伟大，如果今天我们听从祂的声音。今天一天天延长，被称为今天。直到末了，今天和训诲持续；然后真正的今天，天主无尽的日，延续至永远。让我们永远顺服天主圣言的声音。因为今天象征永恒。日子是光的象征；人的光就是圣言，藉着祂我们看见天主。因此，对于相信并服从的人，恩宠将充盈；而对于不信、心里迷误、不认识主道路的人——这些道路若翰吩咐要修直和预备——天主发怒，并威胁他们。\par

而且，确实，从前在旷野中流徙的希伯来人预表性地领受了警告的结局；据说他们因不信没有进入安息，直到跟随了梅瑟的继承人，才亲身体验到（虽然为时已晚），若不信仰耶稣就不能得救。但是，主出于对人的爱，邀请所有人认识真理，为此目的派遣了护慰者。那么，这认识是什么呢？虔诚；按照保禄所言，\BibleQuote{虔诚对一切都有益处，有今生和来生的应许。}\BibleRef{弟前 4:8} 如果永生可以出卖，人们啊，你们打算出多少钱购买？即使有人评估整个\Place{帕克托洛斯河}——传说中金子之河的价钱，也无法估出与救恩相等的价格。\par

但是，不要灰心。如果你愿意，你可以用自己的资本——爱和活泼的信德——购买救恩，虽然它价值无量，这将被视为相称的代价。天主欣然接受这回报；\BibleQuote{因为我们信赖永生的天主，祂是所有人的救主，尤其是信友的救主。}\BibleRef{弟前 4:10}\par

但是其余的人，世上的缠累如同海边的岩石被海藻覆盖一样缠住了他们，他们轻视永生，如同伊萨卡老人，急切地渴望看见的，不是真理，不是天上的祖国，不是真光，而是烟雾。然而，\OriginalTerm{虔诚}{godliness}使人尽可能肖似天主，它指明天主是我们合适的教师，唯有祂能恰当地使人肖似天主。宗徒知道这教导是真正神圣的。他说：\BibleQuote{你，弟茂德，从幼年起就通晓了圣经，这圣经能使你藉着在基督耶稣内的信德，获得得救的智慧。}\BibleRef{弟后 3:15}因为那些书卷实在是神圣的，它们圣化人并使人成神；由这些神圣字母和音节组成的著作或卷册，同一位宗徒随后称之为：\BibleQuote{凡受天主默感所写的圣经，为教训、为督责、为矫正、为教导人学正义，都是有益的，好使天主的人成全，适于行各种善工。}没有人会像被主自己——那爱人者——的话语所感动那样，被任何圣人的劝诫所感动。因为这是祂唯一的工作——人类的救恩。因此祂亲自催促他们走向救恩，呼喊说：\BibleQuote{天国临近了。}\BibleRef{玛 4:17}祂转化那些因恐惧而接近的人。同样，主的宗徒恳求马其顿人，成为神圣声音的解释者，他说：\BibleQuote{主临近了；当心你们不要空手被逮住。}\BibleRef{斐 4:5}但你竟如此缺乏恐惧，或更确切地说，如此缺乏信德，以致不相信主自己，或不相信保罗以基督的身份这样恳求说：\Quote{你们尝尝并看看，基督是不是天主？}信德将引导你进入；经验将教导你；圣经将训练你，因为它说：\BibleQuote{孩子们，来，听我，我要教导你们敬畏上主。}\BibleRef{咏 34:12}然后，对于那些已经相信的人，它简短地补充说：\BibleQuote{谁愿渴慕生命，喜爱享见幸福的日子呢？}\BibleRef{咏 34:13}是我们，我们要说——我们是善的敬拜者，我们热切渴望善的事物。那么，远处的人哪，请听，近处的人哪，请听：圣言没有向任何人隐藏；光是普遍的，它照耀\Quote{所有的人}。没有人在圣言面前是\OriginalTerm{辛梅里安人}{Cimmerian}。让我们奔向救恩，奔向重生；让我们众人赶快，好使我们聚集在一份爱内，按照本质合一的统一；并让我们藉着成为善的，一致地追随合一，追寻\OriginalTerm{善的单一元}{the good Monad}。\par

众多合一，从混杂的声音与分裂中产生神圣的和谐，成为同一交响乐，跟随同一合唱指挥和教师——圣言，达到并安息在同一真理中，呼喊\Quote{阿爸，父}。天主欣然接纳祂子女的这一真实呼声——祂从他们那里收到的初果。\par

\SourceRecord{Translated by William Wilson.；Ante-Nicene Fathers；Vol. 2.；Edited by Alexander Roberts, James Donaldson, and A. Cleveland Coxe.；Buffalo, NY: Christian Literature Publishing Co.,；1885.；<http://www.newadvent.org/fathers/020809.htm>.}

% structure: p0001 source=h1 target=section reason=page-title

\section{劝勉外邦人（第十章）}

\Signature{克肋孟·亚历山大}\par

% structure: p0003 source=h2 target=subsection reason=relative-heading-rank; base=h2

\subsection{第十章 答复外邦人的反对：放弃祖先的习俗是不对的。}

但你们说，推翻从祖先传下来的习俗是不可取的。那么，为何我们仍然不使用最初的食物乳汁，即我们的乳母从我们出生时就使我们习惯的那一种？为何我们增加或减少我们的遗产，而不是保持和我们得到时一模一样？为何我们不再在父母胸前呕吐，或者仍做那些我们在婴儿时期、由母亲哺育时所做而被人嘲笑的事情，而是改正了自己，即使我们没有遇到好的导师？那么，如果纵情欲过度，虽有害且危险，却伴随快感，为何我们在生活行为中不放弃那邪恶的、挑动情欲的、不敬天主的习惯，即令我们的父亲感到痛心，而投奔真理，寻求真正是我们的父亲的那一位，拒绝习俗如同有害的毒品？因为在我所承担的一切任务中，我现在尝试的任务是最崇高的，即向你们证明这疯狂且最可悲的习俗对虔敬是多么敌对。因为如此巨大的恩赐——天主赐予人类的最大恩赐——决不会被憎恨和拒绝，除非你们被习俗冲昏了头脑，然后对我们充耳不闻；就像难以驾驭的马挣脱缰绳、咬住嚼子一样，你们逃离向你们提出的议论，渴望摆脱我们这些试图引导你们生命战车的人，并被你们的愚昧驱使，冲向毁灭的悬崖，将天主圣言视为被诅咒的东西。因此，你们选择的报应，正如\Person{索福克勒斯}所描述的，如下：——\par

\Quote{理智空白，耳朵无用，思想虚妄。}\par

你们不知道，在一切真理中，最真实的是：善良而虔诚的人将获得善报，因为他们高度重视善；反之，恶人将受到应得的惩罚。因为为恶的创始者，刑罚已经预备；所以先知\Person{匝加利亚}警告他说：\BibleQuote{那选择了耶路撒冷的，愿他责斥你；看，这不是从火中抽出来的一根柴吗？}\BibleRef{匝 3:2} 何等愚昧的渴望自愿死亡，竟根植于人心！为何他们奔向这致死的火棍，以使自己被烧，而他们本可以按照天主、而不是按照习俗高尚地生活？因为天主白白赐予生命；但邪恶的习俗，在我们离世后，给罪人带来无用的悔恨和惩罚。\Emph{悲惨的经验，连孩子都知道}迷信毁灭而虔诚拯救。你们中任何人看看那些侍奉偶像的人：头发纠结，身体被肮脏褴褛的衣服玷污；从不亲近浴池，指甲长得惊人，如同野兽；许多人被阉割，他们显示偶像的庙宇实际上是坟墓或监狱。在我看来，他们是在哀悼诸神，而非敬拜，他们的痛苦更值得同情而非虔敬。看到这些，你们仍然盲目，不愿仰望万有的统治者、宇宙的主宰吗？你们不愿逃离那些地牢，逃向从天而降的仁慈吗？因为天主，出于对人的大爱，来帮助人，如母鸟飞向它的一只跌出巢穴的雏鸟；如果一条蛇张口要吞吃小鸟，\Quote{母亲盘旋，发出哀鸣，为她的爱子啼泣}；圣父寻找祂的受造物，治愈他的过犯，追逐蛇，救回幼雏，并激励它飞回巢穴。\par

这样，迷路的狗凭气味追踪主人；摔下骑手的马只要主人吹哨就会来到主人的呼唤。据说：\BibleQuote{牛认识主人，驴认识主人的槽；以色列却不认识我。}\BibleRef{依 1:3} 那么，主呢？祂不记念我们的罪恶；祂仍然怜悯，仍然敦促我们悔改。\par

我还要问你们：你们这些天主的手工品，从祂领受了灵魂，完全属于天主的人，却服从另一位主人，更甚者，服事暴君而非合法的君王——服事邪恶者而非善良者——难道你们不觉得这荒唐吗？因为，以真理之名，哪个有理智的人背向善而依附恶？那么，逃离天主与邪魔为伍的人，是什么人？谁可能成为天主的子女，却宁愿作奴隶？谁能在天国作公民、培育乐园、在天空行走、分享生命树和不朽，并沿着发光的云迹破空而行，像\Person{厄里亚}一样看到救恩的雨，却仍奔向\Person{厄瑞玻斯}？有些人，像虫子在溪流的沼泽和泥泞中打滚，以愚昧无用的享乐为食——猪猡般的人。因为据说猪喜欢泥巴胜过清水，且按\Person{德谟克利特}之说，\Quote{迷恋污秽}。\par

让我们不要作奴隶，也不要成为猪豕一般；而应作为光明的真子女，举目仰望光明，免得主发现我们是假冒的，正如太阳发现鹰隼一般。因此，让我们悔改，从无知走向知识，从愚昧走向智慧，从放纵走向自制，从不义走向正义，从无神走向天主。这是一项高尚的冒险事业：走上通往天主的道路；喜爱正义、追求永生的人，能享受许多其他美善，尤其是天主自己藉依撒意亚所说之事：\BibleQuote{有产业是为那些事奉主的人。} \BibleRef{依 54:17} 这产业是高尚而可欲的：不是金银，不是衣裳，这些都是虫蛀之物，地上的东西也遭盗贼侵犯，盗贼的眼被世俗财富所眩；而是那救恩的宝藏，我们必须奔向它，成为喜爱圣言的人。由此，可赞颂的作为降在我们身上，并与我们同乘真理之翼飞翔。这就是天主永恒的盟约所赋予我们的产业，传递永久的恩宠之赐；因此我们慈爱的圣父——真圣父——不停地劝勉、告诫、训练、爱护我们。因为祂不停地拯救，并指示最好的道路：\Quote{成为正义的，} 主说。\BibleQuote{你们一切干渴的，都来就近水来；没有银钱的，也来，买了喝，不用银钱。} \BibleRef{依 55:1} 祂邀请人到洗礼池、到救恩、到光照，几乎是在呼喊说：我赐给你大地、海洋、天空，还有其中一切的活物，我都白白赐给你。唯愿你，孩子，渴慕你的圣父；天主必白白地显示给你；真理不是买卖之物。祂赐给你空中的飞鸟、水中的游鱼、地上的走兽。这些是圣父为你感恩的享用所创造的。那悖逆之子、灭亡之子、注定做玛门奴隶的人，必须用钱买的，祂却赐给你作为你的产业，甚至给那爱圣父的儿子；为了他，祂仍工作，并且唯独向他应许说：\Quote{地不可永远出卖}，因为它不是注定朽坏的。\Quote{因为全地都是我的}；如果你接受天主，它也是你的。因此，圣经果然向那些相信的人传报喜讯。\Quote{主的圣徒必要承受天主的荣耀和祂的大能。} 什么样的荣耀呢？请你告诉我，蒙福者，是\BibleQuote{眼未曾看见，耳未曾听见，人心也未曾想到的}\BibleRef{格前 2:9}，并且\Quote{他们必在主的王国里永远欢喜！阿们。} 你们，人哪，有天主恩宠的应许；另一方面，也听到了惩罚的警告：主藉这些来拯救，用恐惧和恩宠教训人。我们为何迟延？为何不逃避惩罚？为何不接受白白的恩赐？总之，为何不选择那更好的部分，即天主而非邪恶者，宁取智慧而不拜偶像，以生命换取死亡？\Quote{看，} 祂说，\Quote{我将死亡与生命摆在你们面前。} \BibleRef{申 30:15} 主试验你们，叫你们\Quote{拣选生命}。祂像父亲一样劝你们顺从天主。\Quote{因为，若你们听我，} 祂说，\Quote{且愿意，就必吃地上的美物。} \BibleRef{依 1:19} 这是顺从所伴随的恩宠。\Quote{但若你们不听我，且不愿意，刀剑和烈火必吞灭你们。} 这是悖逆的惩罚。因为主的口——真理的法律，主的话——说了这些。你愿意我作你的好顾问吗？那么，听听吧。我，可能的话，要解释。你们，人哪，在思想美善时，本应提出一个天生的、胜任的见证，即信德，它自己凭自己的资源立即选择最好的，而不是辛苦地探究是否应当追随最好的。因为，容我告诉你们，你们本应怀疑是否应该醉酒，但你们在思考之前就已醉酒；本应怀疑是否应当伤害人，但你们却极其轻易地伤害人。唯一你们拿来质疑的，是是否应当敬拜天主，是否应当追随这位智慧的天主和基督；你们认为这需要深思和怀疑，却不知道什么是配得天主的。你们在我们身上有信心，如同你们在醉酒上有信心，从而成为智慧；你们在我们身上有信心，如同你们在伤害上有信心，从而可以生存。但如果你们承认德行明显的可信，愿意相信它们，来，我要向你们大量陈设关于圣言的劝服材料。但是你们——因为祖传的风俗占据了你们的心智，使你们偏离真理——现在请听真实的情况如下。\par

不要让对这个名字的任何羞耻占据你们，它对人造成极大的伤害，引诱他们离开救恩。那么，让我们公开脱去衣服参加竞赛，在真理的竞技场中高尚地奋斗，圣言是裁判，宇宙的主宰规定竞赛。因为摆在我们面前的并非微不足道的奖赏，即不朽的报酬。因此，如果你们被一些率领不敬之舞的卑贱暴民评头论足，因他们的愚拙和疯狂而驱向同一个坑——即制造偶像和崇拜石头的人——就不要再理会。因为这些人竟敢把人神化——例如马其顿的亚历山大，他们尊他为第十三位神，但巴比伦驳斥了他的狂妄，显示他已死去。因此，我钦佩那位神圣的诡辩家。他名叫泰奥克里托斯。亚历山大死后，泰奥克里托斯拿着人们对神祇所怀的虚妄见解，在同胞面前加以嘲讽，说：\Quote{诸位，只要你们看到神比人死得更快，就请鼓起勇气。}并且，真的，崇拜那些可见的神，以及杂乱的、受生而出的受造物，并依附于它们的人，比那些魔鬼还要悲惨得多。因为天主绝不不义，不像那些魔鬼，而是最高程度地正义；没有什么比我们中一个成为最高程度正义的人更像天主了：——\par

\Quote{上路去，你们所有手工艺人，\newline{}你们崇拜朱庇特怒目而视的女儿，那劳作的女神，\newline{}扇子摆放得当，你们这些愚人}\par

石材的塑造者，崇拜它们的人。让你们的菲迪亚斯、波利克里托斯、普拉克西特列斯和阿佩莱斯也来吧，以及所有从事机械艺术的人，他们自己出自尘土，是尘土的工人。某预言说：\Quote{因为那时，当人信赖肖像时，这里的事务就变得不幸。}让那些更卑微的艺人也来吧——因为我不会停止召唤——。这些人从未制造过有气息的肖像，或用泥土塑造柔软的肉身。谁使骨髓液化？谁使骨头凝固？谁伸展神经？谁扩张血管？谁把血液注入其中？谁铺展皮肤？谁曾制造能观看的眼睛？谁把灵气吹入无生命的形式？谁赐予义德？谁应许不朽？唯有宇宙的创造者；伟大的工匠和圣父造了我们，这样有生命的肖像，即人。但你们的奥林匹斯朱庇特，是肖像的肖像，与真理极不协调，是阿提卡双手的愚拙作品。因为天主的肖像就是祂的圣言，理智的真儿子，神圣的圣言，光之原型光；而圣言的肖像就是真人，即人内在的理智，因此据说人是在\BibleQuote{在天主的肖像和模样内}\BibleRef{创 1:26}，在灵魂的情感中相似神圣的圣言，因此是有理性的；但雕刻成人类形状的雕像，即人可见和尘土所生的那部分的地上肖像，只不过是人性的易朽印记，显然远离真理。那么，那种如此热切于物质的生命，在我看来无非是充满疯狂。而习俗，使你们尝到束缚和不合理的忧虑，被虚妄的意见所滋养；无知，它已证明是人类不合法仪式和骗人表演，以及致命瘟疫和可憎肖像的原因，通过制造许多魔鬼的形状，在所有追随它的人身上烙上持久死亡的印记。那么，接受圣言的水吧；你们这些被污染的，洗濯吧；用真理的水滴洒在你们身上，净化自己脱离习俗。纯洁的人必须升到天堂。如果你看那最普遍于你和他人的，你是人——寻求那位创造了你的；如果你看那独属于你的特权，你是儿子——承认你的圣父。但你仍继续在罪中，沉溺于享乐吗？主将对谁说：\Quote{天国是你们的吗？}你们的，如果你们立志选择天主；你们的，如果你们只相信并顺从这宣告的简短条件；尼尼微人顺从了它，从而避免了他们预期的毁灭，得到了显著的拯救。那么，有人会说：我如何能升到天堂呢？主就是道路；一条窄路，但从天堂而来，实际上窄，却引回天堂，窄，在世受轻视；宽，在天堂受敬拜。\par

然后，那在道理上未受教的人，有愚昧作为其错误的托辞；但至于那已聆听了教诲，却又故意在心里坚持不信的人，他显得越有智慧，他的理智就越会伤害他；因为他自己的理智就是他的控告者，指责他没有选择最好的部分。因为人的本性原本被造就为与天主共融。因此，正如我们不强迫马去犁地，也不强迫牛去狩猎，而是让每种动物去做它本性适合的事；同样，我们指出人与其他受造物不同的独特与突出特征，邀请他——他天生是为默观上天而生，他确实是一株属天的植物——认识天主，劝他为自己准备那足以供给永生的资粮，即虔诚。我们说，如果你是个农夫，就从事农耕；但当你耕种田地时，要认识天主。你，那投身航海的人，航行大海，同时要呼求天上的舵手。你在军中服役时，知识抓住你了吗？听从那命令正当事物的统帅。那么，你们这些被睡梦和醉意压倒的人，要醒过来；稍微睁开你们的眼睛，思考一下你们所崇拜的那些石头意味着什么，以及你们轻率浪费在物质上的支出。你们的资财和产业，你们浪费在愚昧上，就如同把你们的生命抛向死亡，除了这个，你们没有找到虚妄希望的其他结局。你们不仅不能怜悯自己，甚至也不能听从那些同情你们的劝导；你们被邪恶的习俗所奴役，自愿紧抓它直到最后一息，冲向毁灭：\Quote{因为光明来到世界上，世人却爱黑暗甚于光明，}\BibleRef{若 3:19}他们本可扫除那些救恩的障碍——骄傲、财富和恐惧——吟咏着这些诗句：——\par

\Quote{我带着这些丰盛的财富往何处去？\newline{}我自己又往何处徘徊？}\par

那么，你若愿意抛弃这些虚妄的幻影，告别邪恶的习俗，就对虚妄的意见说：——\par

\Quote{说谎的梦，永别了；你们当时什么都不是。}\par

诸位，你们想，提丰的赫尔墨斯、安多基德斯的赫尔墨斯和阿米厄托斯的赫尔墨斯是什么？难道不是显而易见的石头，正如真正的赫尔墨斯本身一样？正如光晕不是神，虹霓也不是神，而是大气和云的状态；同样，一天不是神，一年不是神，由这些组成的时间也不是神；那么太阳和月亮——由它们确定上述事物——也不是神。因此，神智正常的人谁能想象\Quote{纠正}、\Quote{惩罚}、\Quote{公义}和\Quote{报应}是神？因为复仇女神、命运女神、命运本身都不是神，同样政府、荣耀、财富也不是神（画家把后者即普路托斯描绘成瞎子）。但如果你们把谦逊、爱和维纳斯神化，那么就让污名、情欲、美和交媾也随之神化。因此，睡眠和死亡也不能再被合理视为双胞胎神明，它们只是生物身上自然发生的变化；你们也不会恰当地称幸运、命运或命运女神为女神。如果冲突和战争不是神，那么阿瑞斯和厄倪俄也不是神。进一步说，如果闪电、雷霆和雨水不是神，火和水怎能是神？流星和彗星——它们由大气变化产生——怎能是神？谁若称幸运为神，也让他称行动为神。\newline{}那么，如果这些事物以及人手所造的、没有感觉的像都不是神，而关于我们的眷顾显然来自一种神圣能力，那么只剩下承认这一点：那真正是神的那位，只有祂真正存在和存有。但那些对此无感觉的人，就像喝了曼德拉草或其他药的人一样。愿天主赐予你们最终从这沉睡中醒来，认识天主；愿金子、石头、树木、行动、苦难、疾病、恐惧在你们眼中都不再是神。\newline{}因为事实上，\Quote{在肥沃大地上有三万}魔鬼，既非不朽，也非有死；因为它们没有感觉，不足以使其能死，只是木石之物，借着习俗的力量对人实行专制，侮辱并残害生命。\Quote{大地属于主}，经上说，\Quote{及其中的万物。}那么，你们在主恩赐中享乐，为何竟敢无视这至高主宰？\Quote{离开我的地，}主会对你们说，\Quote{不要碰我赐予的水，不可享用我耕种所出产的果实。}为你们的生计回报天主，承认你们的主。你们是天主的受造物。属于祂的东西，怎能公正地被剥夺？因为被剥夺了其所属特性的东西，也就被剥夺了真理。\newline{}因为，效仿尼俄柏的模式，或者更奥秘地说，像古人称为罗特妻子的希伯来女子，你们岂不是变成了一种无感觉的状态？我们听说，这女人因爱索多玛而变成了石头。而那些无神、不虔、心硬、愚昧的人，就是索多玛人。要相信这些话是从天主向你们说的。因为不要以为石头、木块、鸟、蛇是神圣的，而人不是；相反，要把人视为真正神圣的，把野兽和石头当作它们自身。因为人类中有可怜虫，认为天主通过乌鸦和寒鸦发声，而不通过人说话；并且尊崇乌鸦为天主的使者。但天主的人——他不呱噪也不叽喳，而是理性地说话、慈爱地教导——唉，他们竟迫害他；当他邀请他们修养义德时，他们却非人地试图杀死他，既不接纳从上而来的恩宠，也不畏惧惩罚。因为他们不信天主，也不理解祂的大能；祂对人的爱是无法言喻的，祂对恶的憎恨是不可思议的。祂的怒气加剧对罪的惩罚；祂的爱将祝福赐予悔改。失去从天主而来的帮助是极致的悲惨。因此，眼睛的盲目和耳朵的迟钝比恶者的其他折磨更痛苦；因为前者剥夺了他们的天上异象，后者夺走了他们的神圣教导。\newline{}但你们，这样在真理上残缺，心盲耳聋，却不忧伤，不痛苦，不曾渴望看见天堂和天堂的创造者，也不曾通过选择救恩而寻求听见宇宙的创造者并向祂学习；因为对于决心认识天主的人，没有任何障碍。无论无子女、贫穷、卑微、匮乏，都不能阻止热切追求认识天主的人；也没有人凭铜铁为自己征服了真智慧而选择放弃它，因为它远远胜于一切。基督能在各处拯救。因为那满怀义德热情和钦佩的人，是那无需万物的独一者的爱人，他自己所需极少，只在天主和自己那里积存了福乐，在那里没有蛾子、\BibleRef{玛 6:20}强盗或海盗，只有永恒的赐恩者。\newline{}因此，你们恰被比作那些对弄蛇者掩耳的蛇。因为\Quote{他们的心思}，经上说，\Quote{就像蛇，像聋的毒蛇，它堵住耳朵，不愿听弄蛇者的声音。}但请允许自己感受圣洁的迷人旋律，接受我们温和的话语，拒绝致命的毒药，这样你们就有可能尽可能脱去毁灭，如同他们脱去了老年。听我说，不要堵住耳朵，不要堵塞听觉的通道，要留心所的话。永生的良药何其美妙！停止你们匍匐爬行的动作。\Quote{因为主的仇敌}，经上说，\Quote{必舔尘土。}抬起你们的眼睛从地上看天空，仰望天堂，欣赏那景象，别再伸长头盯着义人的脚跟，阻碍真理之路。要明智而无害。也许主会赐给你们纯朴的翅膀（因为祂已决定给土生者翅膀），使你们离开洞穴，居住在天上。只要我们全心全意悔改，就能全心全意容纳天主。\Quote{你们所有聚会的百姓，要信赖祂，在祂面前倾吐你们的心。}祂对那些刚离弃邪恶的人说：\Quote{祂怜悯他们，以义德充满他们。}相信那既是人又是天主的祂；相信吧，人啊。相信吧，人啊，那活着、受苦并被敬拜的天主。相信吧，你们这些奴隶，那死去的祂；相信吧，所有人类，那独一的全人类的天主。相信吧，并以救恩作为你们的奖赏。寻求天主，你们的灵魂必得生存。寻求天主的人是在关心自己的救恩。你们找到了天主吗？那么你们就有了生命。让我们寻求，为使我们得以生活。寻求的奖赏是与天主同在的生命。\Quote{愿所有寻求祢的人都因祢欢欣鼓舞，愿他们不断说：愿天主受显扬。}一首高贵的赞美天主的诗歌，是一个在义德中建立的不朽之人，真理的话语铭刻在他身上。因为除了在智慧的灵魂里，你能在哪里写下真理？爱在哪里？敬畏在哪里？温柔在哪里？那些印有这些神圣特质的人，我认为，应当视智慧为美丽的港口，无论他们转向何种人生境遇；同样也视它为救恩的平静港湾：智慧使那些投靠了父的人成为子女的好父亲；使那些认识了子的人成为儿子的好父母；使那些记念新郎的人成为妻子的好丈夫；使那些从彻底奴役中得赎的人成为仆人的好主人。\newline{}哦，陷入错误的人啊，野兽比你们幸福得多！谁生活在无知中如同你们，却不冒充真理。他们中没有谄媚的族类。鱼没有迷信；鸟不敬拜任何一个偶像；它们只仰慕天堂，因为它们缺乏理性，无法认识天主。你们难道不羞愧，在如此漫长的生命中活在不虔之中，使自己比无理性的动物更无理性吗？你们曾是孩童，然后少年，然后青年，然后成人，却从未成为善人。如果你们尊重老年，现在到了生命的日落时分，就变明智吧；即使在生命结束时，也要获得对天主的认识，使生命的终结成为救恩的开端。你们在迷信中变老了；要像年轻人一样进入虔敬的实践。天主视你们为无辜的孩子。\newline{}那么，让雅典人遵守梭伦的法律，阿耳戈斯人遵守福洛纽斯的法律，斯巴达人遵守吕库古的法律；但如果你们将自己登记为天主的子民，天堂就是你们的国家，天主就是你们的立法者。法律是什么？\Quote{不可杀人；不可奸淫；不可诱拐男童；不可偷盗；不可作假见证；要爱主你的天主。}\BibleRef{出 20:13}而补充这些的是那些铭刻在人心的理性法律和圣洁话语：\Quote{要爱你的邻人如同你自己；打你脸颊的，把另一面也转给他；}\BibleRef{路 6:29}\Quote{不可贪恋，因为单凭贪恋你就已犯了奸淫。}\BibleRef{玛 5:28}那么，人们从起初就不愿贪恋禁止之事，比起得到所欲之物，岂非好得多！但你们不能忍受救恩的严肃；而我们喜爱甜物，因它们带来的愉悦而珍视它们；另一方面，那些口感苦涩之物却有疗效和医治作用，良药的苦口增强了胃弱之人的力量——这样习俗令人愉悦和瘙痒；但习俗将人推入深渊，真理却引向天堂。起初它严厉，却是青春的好保姆；它既是女仆主妇共处的端庄之所，又是智慧议事厅。它并不难接近或不可达到，而是近在我们的家里；正如充满智慧的梅瑟所说，提及它时，它在我们构成的三部分中有住所——在手、口和心：这是真理的恰当象征，为众人通过意愿、行动和言语这三者所拥抱。\newline{}不要害怕眼前升起的众多悦人之物会使你离开智慧。你会自发地超越习俗的轻浮，如同男孩长大成人后抛弃玩具。因为以无与伦比的速度和容易接近的仁慈，神圣的大能将光辉洒在地上，使宇宙充满了救恩的种子。因为并非没有神圣的眷顾，如此伟大的工程在主内于如此短的时间完成；祂在外表上被轻视，实际上却被敬拜，祂是赎罪者、救主、仁慈者、圣言、真正最显明的神性、与宇宙之主同等的那位；因为祂是祂的子，圣言在天主内，在最初被宣讲时并非所有人都怀疑，当祂取了人的身份，在肉身中塑造自己，上演人类救恩的戏剧时，也并非完全不被认识：因为祂是真正的勇士，与受造物并肩作战。并且以最快速度传递给人类，从父的旨意中升起比太阳更快，以最完全的方式使天主照耀我们。祂从何处来、祂是什么，由祂所教导和展示的表明，祂彰显自己为盟约的使者、和好者、我们的救主、圣言、生命之泉、和平的赐予者，遍及全地面；经由祂，可以说，宇宙已经成为祝福的海洋。\par

\SourceRecord{Translated by William Wilson.；Ante-Nicene Fathers；Vol. 2.；Edited by Alexander Roberts, James Donaldson, and A. Cleveland Coxe.；Buffalo, NY: Christian Literature Publishing Co.,；1885.；<http://www.newadvent.org/fathers/020810.htm>.}

% structure: p0001 source=h1 target=section reason=page-title

\section{\WorkTitle{劝勉外邦人（第十一章）}}

\Signature{亚历山大城的克肋孟著}\par

% structure: p0003 source=h2 target=subsection reason=relative-heading-rank; base=h2

\subsection{第十一章 基督来临为人带来何等伟大的恩惠}

你若有兴趣，请稍加默观天主的慈惠。第一人在乐园中自由嬉戏，因他是天主的儿女；但他一旦屈服于享乐——因为蛇寓意性地表示腹行于地的享乐，即地上邪恶欲火之燃料——便如孩童般被情欲引诱，在悖逆中衰老；因违抗父命而侮慢了天主。享乐的影响即如此。原本因纯朴而自由的人，竟被罪捆绑。主于是愿意释放他脱离锁链，披上肉身——啊，神圣的奥迹！——战胜了蛇，奴役了死亡的暴君；而最令人惊叹的是，那被享乐欺骗、被腐朽牢牢束缚的人，双手被解开，得了释放。啊，神秘的奇事！主被贬抑，人却升起；那从乐园坠落的人，因顺服而得着比乐园更大的赏报，即天堂本身。因此，既然圣言亲自从天上来到我们中间，我认为我们不再需要前往雅典、希腊其余地方及伊奥尼亚去寻求人的学问。因为如果我们有那位以祂神圣德能在创造、救恩、恩惠、立法、预言、教导中充满宇宙的作为我们的导师，我们便有了那一切教导之源的导师；而整个世界，连同雅典和希腊，已成了圣言的疆域。你们既相信那诗意的寓言，称克里特的米诺斯是宙斯的密友，就不会不信我们这些成为天主的门徒的人已获得了唯一真正的智慧；而哲学首领们仅猜测到的，基督的门徒却已领悟并宣扬。那独一的整个基督并不分裂：\Quote{没有外邦人、犹太人、希腊人，没有男或女，只有一个新人}，由天主的圣神转化而成。此外，其他的劝谕和诫命无关紧要，只关乎具体事物——例如，可否结婚、参与公共事务、生养儿女；但唯一普遍性的、贯穿整个生存过程、时时处处皆引向最高目标即生命的命令，乃是虔诚——我们为得永生所需要的一切，都在于按此生活。然而哲学，如古人所言，是\Quote{一种长久劝勉，追求智慧永恒的爱}；而主的诫命则远照，\Quote{照亮眼睛}。接受基督，接受视觉，接受你的光明，\par

\Quote{为使你清楚认识天主和人。}\par

\Quote{赐我们光明的圣言是甘饴的，贵于黄金和宝石；比蜜和蜂房更可羡慕。}因为它怎能不是可羡慕的呢？它已用光照亮了那曾埋在黑暗中的心灵，并赋予了灵魂\Quote{光明之眼}敏锐的视力。因为，正如倘若没有太阳，即使有其他天上的光体，黑暗仍会笼罩宇宙；同样，如果我们未曾认识圣言、未蒙祂光照，我们与那些在黑暗中养肥、被喂养以待宰杀的飞鸟毫无区别。那么，让我们接受光明，好让我们接受天主；让我们接受光明，成为主的门徒。这也是祂已向父所应许的：\Quote{我要向我的弟兄们宣告祢的名；在教会中我要赞美祢。}赞美并向我宣告祢的父天主；祢的话语拯救；祢的赞歌教导我：此前我一直在迷途中徘徊，寻求天主。但因祢领我到光明，主啊，我借着祢找到了天主，从祢接受了父，我成了\Quote{祢的同承继人}\BibleRef{罗 8:17}，因为祢\Quote{不以我为弟兄为耻}\BibleRef{希 2:11}。那么，让我们抛弃，让我们抛弃对真理的遗忘，即无知；除去那如同视障的阻碍的黑暗，让我们默观唯一真天主，首先以这赞美之歌高呼：万福，光明！因为在我们这些埋在黑暗中、被困于死影里的人身上，从天上照耀了比太阳更纯净、比现世生命更甘饴的光明。那光明就是永生；凡分沾它的便活着。但黑夜惧怕光明，恐惧地躲藏，给主的日子让位。不眠的光明如今普照一切，西方已相信东方。因为这是新创造的目的。因为\Quote{正义的太阳}，祂驾着战车巡行一切，同样遍及全人类，如同\Quote{祂的父，使祂的太阳升起照所有人}，并将真理的甘露降在他们身上。祂将日落变为日出，借着十字架使死亡成为生命；祂将人从毁灭中夺出，高举至诸天，把必朽的移植为不朽，把大地迁至天堂——祂，天主的农夫，\par

\Quote{指出吉兆，唤醒万民\newline{}行善，提醒他们真正的食粮；}\par

将父那真正伟大、神圣、不可剥夺的遗产赐予我们，以天上的教导使人成神，将祂的律法放在我们心里，写在我们的心上。祂刻下什么律法？\Quote{所有人都要认识天主，从小到大的都认识}；并且，\Quote{我要怜悯他们}，天主说，\Quote{不再记念他们的罪}\BibleRef{希 8:10}。让我们接受生命的律法，让我们顺从天主的劝戒；让我们认识祂，好使祂施恩。虽然天主一无所缺，让我们以感恩的心和虔诚作为我们在这地上居所的租金，向祂献上感恩的回报。\par

\Quote{用金换铜，\newline{}用一百头牛的代价换九头牛；}\par

就是说，为了你的小信德，祂赐给你如此广阔的土地耕种，赐给你水可饮也可航行，赐给你空气呼吸，赐给你火为你效力，赐给你世界居住；祂还允许你率领一群殖民从这里前往天堂；藉着祂手的这些重要化工和如此众多的恩惠，祂酬报了你小小的信德。可是，那些信靠\OriginalTerm{招魂术士}{necromancers}的人，从他们那里得到\OriginalTerm{护身符}{amulets}和\OriginalTerm{咒符}{charms}，以驱避邪恶；难道你们就不肯让天上的圣言，救主，像护身符一样系在你们身上，藉着信赖天主的咒符，从作为心灵疾病的私欲偏情中被解救，从罪中被拯救出来吗？——因为罪是永死。确实，你们完全迟钝、盲目，像鼹鼠一样，除了吃无所事事，在黑暗中度过一生，被腐朽包围。但真理呼喊说：\Quote{光要从黑暗中照耀出来。}那么，让光照耀在人隐藏的部分，即心里；让知识的曙光升起，以揭示和照亮隐藏的内在之人，那光的门徒，基督的亲密朋友和共同继承者；尤其如今我们已认识了良善圣父最珍贵、最可敬的名号，祂对虔诚而良善的儿女给予柔和的训诲，并吩咐那对祂儿女有益的事。那服从祂的人在一切事上都有益处，追随天主，服从圣父，经过漂泊认识祂，爱天主，爱近人，履行诫命，寻求奖赏，索取恩许。但是天主固定不变的目的始终是要拯救人类的羊群：为此，良善的天主派遣了良善的牧人。圣言展开了真理，向人显示了救恩的高度，好叫他们或悔改而得救，或拒绝服从而受审判。这是正义的宣告：对服从者，是喜讯；对不服从者，是审判。响亮的号角吹响时，召集士兵，宣告战争。基督岂不向地极吹奏和平的曲调，聚集祂自己的士兵，即和平的士兵吗？的确，藉着祂的血和圣言，祂聚集了和平的无血大军，将天国赐给了他们。基督的号角就是祂的福音。祂已吹响，我们都听到了。\BibleQuote{让我们穿上和平的铠甲，佩上正义的胸甲，拿起信德的盾牌，用救恩的头盔束住额头；并拿起圣神的剑，即天主的圣言，} \BibleRef{弗 6:14} 来磨快吧。宗徒怀着和平的精神如此命令。这些是我们无敌的武器：以这些武装，让我们面对恶者；\Quote{恶者的火箭}让我们用蘸水的剑尖熄灭，好使我们因圣言受洗，对所领受的恩惠献上感恩，并藉着神圣圣言荣耀天主。\Quote{因为当你还在说话时，}经上说，\Quote{祂会说：看，我就在你身边。} \BibleRef{依 58:9}哦，这神圣而蒙福的能力，天主藉此与人共融！那么，远为更好的是立刻成为万有最佳者的效仿者和仆人；因为只有藉着神圣的侍奉，人才能效仿天主，而只有藉着效仿祂，才能侍奉和朝拜祂。天上的、真正神圣的爱就这样来到人身上：当在灵魂本身之中，由神圣圣言点燃的真正良善的火花，能够爆发成火焰；而最重要的是，救恩与真诚的意愿并行——选择和生命，可以说是搭配在一起的。因此，这真理的劝勉独自像我们最忠实的朋友一样，陪伴我们直到最后一息，并且对于灵魂圆满完美的精神，是我们升天时慈祥的伴侣。那么，我给你的劝勉是什么呢？我劝你获得救恩。这是基督所愿。一言以蔽之，祂白白赐予你生命。祂是谁？简要地学习。真理的圣言，不朽的圣言，祂藉着将人带回真理而重生人——催迫向救恩的刺棒——祂驱逐毁灭、追赶死亡——祂在人心中建筑天主的殿宇，好使天主住到人里面。洁净这殿；将快乐和娱乐抛给风和火，如同凋谢的花；但要明智地栽培自制的果实，将你自己呈献给天主作为初熟之果，好使不仅有人工，也有天主的恩宠；两者都是必要的，好使基督的朋友得以堪当天国，并被算为配得天国。\par

\SourceRecord{Translated by William Wilson.；Ante-Nicene Fathers；Vol. 2.；Edited by Alexander Roberts, James Donaldson, and A. Cleveland Coxe.；Buffalo, NY: Christian Literature Publishing Co.,；1885.；<http://www.newadvent.org/fathers/020811.htm>.}

% structure: p0001 source=h1 target=section reason=page-title

\section{劝勉外邦人（第十二章）}

\Signature{克莱孟·亚历山大}\par

% structure: p0003 source=h2 target=subsection reason=relative-heading-rank; base=h2

\subsection{第十二章 劝勉放弃旧谬误，听从基督的训诲。}

让我们像躲避危险的海岬、或可怕的\OriginalTerm{卡律布狄斯}{Charybdis}、或神话中的\OriginalTerm{塞壬}{sirens}那样，避开习俗。它窒息人，使人远离真理，引人离开生命：习俗是网罗、深渊、陷阱、恶毒的簸箕。\par

\Quote{把船驶过那烟雾和巨浪。}\par

让我们避开，同船的水手们，让我们避开这道巨浪；它喷出火焰：它是一个邪恶的岛屿，堆满白骨和尸体，岛上有一位美丽的妓女，享乐，用音乐取悦众人之耳。\par

\Quote{来吧，声名远扬的\OriginalTerm{尤利西斯}{Ulysses}，\OriginalTerm{阿开亚人}{Achæans}的伟大荣光；\newline{}停泊船只，好让你聆听更神圣的声音。}\par

她赞美你，啊水手，称你为卓越者，并试图把希腊人的荣耀归给自己。让她在死人身上觅食吧；一位天上的神灵帮助你：越过享乐，她诱惑人。\par

\Quote{不要让拖着长裙的女子欺骗你的感官，\newline{}用她谄媚的闲聊寻求你的伤害。}\par

驶过那歌声；它带来死亡。只要运用你的意志，你就已经克服了毁灭；被绑在十字架的木头上，你将免于败坏：天主圣言将成为你的舵手，圣神将把你锚定于天堂的海港。那时你将看见我的天主，并接受神圣奥秘的入门，享受那些在天上为我所保留的事物，\BibleQuote{耳朵未曾听见，人心也未曾想到。} \BibleRef{格前 2:9}\par

\Quote{实在我看似看见两个太阳，\newline{}和一双\OriginalTerm{底比斯}{Thebes}，}\par

一个在敬拜偶像中疯狂、因纯粹无知而酩酊大醉的人这样说道。我怜悯他癫狂的醉酒，并这样癫狂地邀请他来到救恩的清醒中；因为主欢迎罪人的悔改，而不是他的死亡。\par

来，啊狂人，不要倚靠\OriginalTerm{松果杖}{thyrsus}，不要戴着常春藤冠；扔掉\OriginalTerm{高冠}{mitre}，扔掉\OriginalTerm{鹿皮}{fawn-skin}；恢复理智。我要向你展示圣言和圣言的奥秘，按你的方式加以解释。这是天主所爱的山，不像\OriginalTerm{基塞隆山}{Cithæron}那样成为悲剧的主题，而是奉献给真理的戏剧——一座清醒的山，覆盖着纯洁的森林；在那里狂欢的不是\OriginalTerm{塞墨勒}{Semele}的姐妹\OriginalTerm{迈那得斯}{Mænades}——塞墨勒被雷击，在入门仪式中实行不洁的肉块分割——而是天主的女儿们，美丽的羔羊，她们庆祝圣言的神圣礼仪，跳着清醒的轮舞。义人是合唱队；音乐是宇宙之王的颂歌。少女们弹奏竖琴，天使们赞美，先知们发言；音乐声响起，他们奔跑追逐欢庆的队伍；被召叫的人急忙赶来，热切渴望领受圣父。\par

来，啊老人，离开底比斯，从你身上抛弃占卜和酒神狂热，允许自己被引领到真理。我给你（十字架的）杖作为依靠。快，\OriginalTerm{提瑞西阿斯}{Tiresias}；相信，你就会看见。基督，盲人借祂恢复视力，将为你照耀比太阳更亮的光；黑夜将从你面前逃走，火将畏惧，死亡将消失；你，老人，未见底比斯，却将看见诸天。啊真正神圣的奥秘！啊无瑕之光！我的道路由火炬照亮，我仰观诸天和天主；我在入门时成为圣洁。主是\OriginalTerm{宣教者}{hierophant}，并给入门者盖印同时光照他，并将相信的人呈献给圣父，使他永远被保守。这就是我奥秘的沉思。如果你愿意，也请入门；你将与天使们一同围绕未受生、不可毁灭的唯一真天主——天主圣言，与我们一同高唱赞歌。这位耶稣——永远的、唯一真天主的唯一大司祭，也是祂的圣父的大司祭——为人类祈祷并劝勉。\par

\Quote{听，你们无数民族，更确切地说，凡是被赋予理性的人，无论野蛮人还是希腊人。我呼唤全人类，我是他们的创造者，藉圣父的旨意。来就我，好使你们被安置在唯一天主和唯一天主圣言之下的合宜等级中，并且不仅因拥有理性而比无理性之物优越；因为对你们所有可朽者，我赐予不朽的享受。因为我愿意，我愿意将这恩宠赐予你们，将不朽的完美恩赐给予你们；我将圣言和天主的认识，即我完全的自体，赐予你们。这就是我，这就是天主旨意，这就是交响，这就是圣父的和谐，这就是圣子，这就是基督，这就是天主圣言，上主的手臂，宇宙的能力，圣父的旨意；这些事在古代有预像，但并不完全。我渴望按原初的样式恢复你们，使你们也肖似我。我用信德之油傅抹你们，使你们脱去腐朽，并向你们展示赤裸的正义形体，使你们由此升到天主面前。凡劳苦和负重担的，都到我这里来，我要使你们安息。你们背起我的轭，跟我学，因为我是良善心谦的：这样你们必要找得你们灵魂的安息，因为我的轭是柔和的，我的担子是轻省的。}\par

我们这些热爱天主、相似天主的圣言肖像的同胞啊，让我们快跑，让我们奔跑吧。让我们快跑，让我们奔跑，让我们负起祂的轭，让我们接受引领我们进入永生的、人类的好车夫。让我们爱基督。祂曾带着母驴和小驴驹；祂将人类的队伍与天主联结，驾驭祂的战车奔向永生，祂现在驱车进入天堂，清楚急促地实现了祂先前骑着驴进入耶路撒冷所预表的事。对圣父而言，最美丽的景象是永恒圣子戴着胜利的冠冕。那么，让我们追求善，让我们成为热爱天主的人，获得一切事物中最伟大、不可伤害的——天主和生命。我们的助佑者是圣言；让我们信赖祂；永远不要让我们被对金银和荣耀的渴望所困扰，如同对真理圣言本身的渴望那样。因为如果我们轻视最有价值的事物，却将明显的罪恶、极度的不敬——愚昧、无知、轻率、偶像崇拜——奉为最高，这绝不会，绝不会取悦天主本身。因为哲学家的弟子们不无理由地认为，愚人在所有行为中都犯有亵渎和不敬；他们将无知本身描述为一种疯狂，断言群众不过是疯子。因此（圣言会说）毫无疑问，清醒比疯狂更好；但我们必须用牙齿紧紧咬住真理，全力跟随天主，运用智慧，视万物如同祂的一样；此外，得知我们是他最卓越的财产，让我们将自己交付给天主，爱主天主，将这事视为我们一生的职分。如果朋友间的财物被视作公有，而人是天主的朋友——因为他藉着圣言的中介成为天主的朋友——那么相应地，万物都成为人的，因为万物属于天主，是朋友双方——天主和人——的共同财产。\par

是时候了，我们应当说：唯有虔诚基督徒才富有、智慧、出身高贵，并这样称呼和相信他是天主的肖像和模样，因\Person{耶稣基督}而成正义、圣洁、智慧，并已如此相似天主。因此先知指出了这恩宠，他说：\BibleQuote{我曾说你们是神，都是至高者的儿子。} 对于我们，是的，我们，祂收纳了我们，并愿意被称为惟独我们（不是不信者）的圣父。这就是我们这些基督随从者的处境。\par

\Quote{人的意愿如何，言语也如何；\newline{}言语如何，行为也如何；\newline{}行为如何，生命也如何。}\par

认识基督的人，他们的整个生命是美好的。\par

我想，话已足够。虽然出于对人的爱，我本可以继续倾吐我从天主领受的，好劝勉你们追求最伟大的福分——救恩。因为关于无尽生命的言谈，不容易讲到揭示的尽头。你们仍有这个结论要作：选择哪一样对你们最有益——审判还是恩宠。因为我不认为对两者孰优尚有疑问；也不允许将生命与毁灭相提并论。\par

\SourceRecord{Translated by William Wilson.；Ante-Nicene Fathers；Vol. 2.；Edited by Alexander Roberts, James Donaldson, and A. Cleveland Coxe.；Buffalo, NY: Christian Literature Publishing Co.,；1885.；<http://www.newadvent.org/fathers/020812.htm>.}
